\documentclass[12pt]{article}
\usepackage[margin=1in]{geometry}
\usepackage{amsmath,amssymb,amsthm}
\usepackage{mathrsfs}
\usepackage{graphicx}
\usepackage{hyperref}
\usepackage{booktabs}
\usepackage{listings}
\usepackage{enumitem}
\usepackage{array}

\theoremstyle{plain}
\newtheorem{theorem}{Theorem}[section]
\newtheorem{proposition}[theorem]{Proposition}
\newtheorem{lemma}[theorem]{Lemma}
\newtheorem{corollary}[theorem]{Corollary}
\theoremstyle{definition}
\newtheorem{definition}[theorem]{Definition}
\newtheorem{axiom}{Axiom}
\theoremstyle{remark}
\newtheorem{remark}[theorem]{Remark}

\lstset{basicstyle=\ttfamily\footnotesize, breaklines=true, frame=single}

\title{Stratified Geometric Ethics: Mathematical Foundations for
Verifiable Moral Reasoning in Autonomous Systems}
\author{Andrew Bond\\Department of Computer Engineering\\San Jos\'e State University}
\date{March 2026}

\begin{document}
\maketitle

\begin{abstract}
We present \emph{Stratified Geometric Ethics} (SGE), a mathematical
framework providing rigorous foundations for deterministic, verifiable
ethical reasoning in autonomous systems. SGE addresses fundamental
limitations of prior geometric approaches by modeling the space of
ethically relevant configurations as a \emph{stratified space}---a union
of smooth manifolds of varying dimensions connected along boundary
strata---rather than a smooth manifold. This structure captures moral
discontinuities, incommensurable values, and genuine ethical dilemmas
that smooth models cannot represent.

We make five principal contributions. First, we show that stratified
spaces are a \emph{natural minimal candidate} among standard geometric
structures for representing ethical phenomena including threshold
effects, lexical priorities, and moral dilemmas (Theorem 2.3). Second,
we establish a \emph{representation theorem} (Theorem 4.3)
characterizing all satisfaction functionals satisfying five explicit
axioms---including a novel locality axiom and scale-normalization
assumption---with a complete proof. Third, we prove \emph{finite
approximation theorems} (Theorems 3.9--3.11) showing that any decision
problem on a compact stratified space reduces to a finite graph problem
with explicit error bounds. Fourth, we establish \emph{decidability
results} (Theorem 6.4) for the quantifier-free, non-temporal fragment of
our ethical specification language via o-minimal structures, with
temporal properties handled by standard model checking on finite
approximations. Fifth, we derive \emph{sample complexity bounds}
(Theorems 7.1--7.3) for learning ethical content from data.

We introduce the Bond Invariance Principle (BIP), which requires ethical
judgments to depend only on morally relevant relationships, not
arbitrary representation choices, and prove that our axiom system
implies BIP.

SGE serves as the theoretical foundation for the DEME architecture,
providing mathematical justification for design choices including
multi-dimensional moral vector spaces, governance profiles with veto
regions and lexical priorities, and layered enforcement architectures.
Since initial development, a comprehensive reference implementation
(196 commits, 15 development sprints) has realized all core theoretical
constructs as rank-6 tensors with hardware-accelerated backends,
and the Bond Invariance Principle has been empirically validated across
11 languages (80.0\% F1, 11.1$\times$ structural-to-surface ratio,
\(p = 0.023\)).
\end{abstract}

\noindent\textbf{Keywords:} stratified geometric ethics, Whitney stratification, gauge invariance, moral reasoning, o-minimal structures, formal verification, AI alignment, autonomous systems

\section{1 Problem statement and
approach}\label{problem-statement-and-approach}

The deployment of AI systems in safety-critical domains---healthcare,
autonomous vehicles, criminal justice, financial markets---creates an
urgent need for frameworks that make ethical reasoning explicit,
deterministic, and formally verifiable. Current approaches fall into two
categories: \emph{principle-based} approaches that articulate high-level
commitments without computational implementations, and
\emph{learning-based} approaches that infer ethical behavior from
feedback without formal guarantees. Neither satisfies the requirements
of safety-critical deployment where lives depend on ethical decisions
made at machine timescales.

Recent work has proposed differential-geometric frameworks for ethics,
modeling configurations as points on smooth manifolds and ethical
quantities as tensor fields {[}1{]}. While mathematically elegant, these
approaches face five fundamental limitations:

\textbf{The Content Problem.} Existing frameworks assume obligations and
interests are "given" without specifying how these tensors derive from
ethical sources---constitutional principles, stakeholder preferences, or
deliberative procedures.

\textbf{The Smoothness Problem.} Real ethical reasoning involves
discrete choices (trolley problems), incommensurable values (life vs.
property), and threshold effects (killing vs. letting die). Smooth
manifolds cannot represent these discontinuities.

\textbf{The Linearity Problem.} Bilinear satisfaction functionals S(x) =
O\textsuperscript{μ}(x)I\textsubscript{μ}(x) cannot capture threshold
effects, lexical priorities, or diminishing returns---phenomena central
to moral reasoning.

\textbf{The Tractability Problem.} No existing geometric ethics
framework provides complexity analysis establishing whether ethical
computation can be performed in realtime.

\textbf{The Verification Problem.} No framework enables
machine-checkable proofs of ethical constraint satisfaction.

This paper introduces \emph{Stratified Geometric Ethics} (SGE), a
framework that addresses all five limitations while preserving the
coordinate-invariance that makes geometric approaches valuable. The
companion paper {[}2{]} instantiates SGE in the DEME 2.0 architecture
for real-time ethical governance in autonomous systems (now upgraded to DEME 3.0 with tensorial ethics support {[}16{]}).

This paper argues that a central failure mode in machine ethics is not
merely bias in data, but \textbf{representation dependence}: conclusions
that shift when only the description changes. We formalize a bond-based
invariance principle and derive the constrained class of admissible
evaluators it entails.

\subsection{1.1 Relationship to DEME}\label{relationship-to-deme}

SGE was developed through a process of reflective equilibrium with DEME,
a computational architecture for real-time ethical governance
described in the companion paper {[}2{]}. DEME commits to concrete
engineering choices: multi-dimensional moral vector spaces, governance
profiles with veto regions and lexical priorities, a three-layer
architecture (strategic, tactical, reflex), and hardware-resident ethics
modules. Since initial development, the architecture has been extended
to V3 with rank-6 tensor operations, multi-agent coordination, temporal
dynamics, and uncertainty quantification (Section 10). SGE provides the
mathematical justification for these choices:

\begin{quote}
• DEME\textquotesingle s \emph{moral vector space} M ⊆
ℝ\textsuperscript{k} is justified by Theorem 2.3, which shows stratified
spaces are natural minimal candidates for ethical phenomena.

• DEME\textquotesingle s \emph{governance profiles} with veto regions
and scalarization are justified by Theorem 4.3\textquotesingle s
representation of satisfaction operators.

• DEME\textquotesingle s \emph{real-time enforcement} is justified by
the finite approximation theorems (3.9--3.11) and complexity bounds
(Section 5).

• DEME\textquotesingle s \emph{verification layer} is justified by the
decidability results (Theorem 6.4) combined with standard model
checking.
\end{quote}

\section{2 Why Stratified Spaces}\label{why-stratified-spaces}

Before developing the machinery of SGE, we establish that stratified
spaces are a \emph{natural minimal candidate} among standard geometric
structures for capturing ethical phenomena. We do not claim a fully
formal minimality theorem over all conceivable geometric
structures---such a claim would require defining "simpler" in a way that
excludes exotic alternatives. Instead, we show that standard
alternatives (smooth manifolds, manifolds with corners, cell complexes)
each fail to represent at least one essential ethical phenomenon, and
that stratified spaces succeed at all of them.

\subsection{2.1 Ethical Phenomena Requiring Geometric
Representation}\label{ethical-phenomena-requiring-geometric-representation}

We identify four categories of ethical phenomena that any adequate
geometric framework must represent:

\textbf{E1. Discrete Choices.} Many ethical decisions present finitely
many options with no meaningful intermediate. In trolley-type scenarios,
one must choose path A or path B; there is no "path 0.5A + 0.5B."

\textbf{E2. Incommensurable Values.} Some moral dimensions cannot be
traded off at any finite rate. The value of a human life is not
equivalent to any amount of property damage; procedural rights cannot be
"compensated" by increased welfare.

\textbf{E3. Threshold Effects.} Crossing certain boundaries has
discontinuous moral significance. The distinction between killing and
letting die, between lying and remaining silent, between 0.99 and 1.00
on a consent scale---these involve discrete moral transitions.

\textbf{E4. Genuine Dilemmas.} Some situations admit no ethically
satisfactory resolution. Both available options involve genuine moral
loss ("moral residue"), and this property must be represented, not
smoothed away.

\subsection{2.2 Insufficiency of Smooth
Manifolds}\label{insufficiency-of-smooth-manifolds}

\textbf{Proposition 2.1.} \emph{Let M be a smooth connected manifold.
Then:}

\begin{quote}
\emph{(i) Any two points in M can be connected by a smooth path
(violates E1);}

\emph{(ii) Any smooth function f : M → ℝ has continuous level sets
(violates E3);}

\emph{(iii) For any smooth Riemannian metric g on M, all directions at
any point have finite, comparable costs (violates E2).}
\end{quote}

\emph{Proof.} (i) follows from path-connectedness of connected
manifolds. (ii) follows from continuity of smooth functions. (iii)
follows from positive-definiteness of Riemannian metrics: g(v,v)
\textgreater{} 0 for all nonzero v, so all directions have positive
finite cost, hence are comparable. ■

\subsection{2.3 Stratified Spaces as Natural Minimal
Candidates}\label{stratified-spaces-as-natural-minimal-candidates}

\textbf{Definition 2.2} (Stratified Space)\textbf{.} \emph{A stratified
space is a triple (M,
\{M}\textsubscript{i}\emph{\}}\textsubscript{i∈I}\emph{, ⪯) where M is a
paracompact Hausdorff space, \{M}\textsubscript{i}\emph{\} is a locally
finite partition into connected smooth manifolds (strata), and ⪯ is a
partial order on I such that M}\textsubscript{i} \emph{∩
cl(M}\textsubscript{j}\emph{) ≠ ∅ implies i ⪯ j (frontier condition). We
require Whitney\textquotesingle s condition (B) for regularity.}

\textbf{Theorem 2.3} (Stratified Spaces Represent E1--E4)\textbf{.}
\emph{Stratified spaces can represent all four ethical phenomena:}

\begin{quote}
\emph{(i) Discrete choices are modeled by 0-dimensional strata
representing decision endpoints;}

\emph{(ii) Incommensurable values are modeled by stratified metrics with
arbitrarily large cost (in a singular limit);}

\emph{(iii) Threshold effects are modeled by stratum boundaries where
satisfaction functions have discontinuities;}

\emph{(iv) Genuine dilemmas are modeled by singular strata from which
all exits incur positive moral cost.}
\end{quote}

\emph{Moreover, among standard geometric structures---smooth manifolds,
manifolds with corners, and cell complexes without smooth
structure---each fails at least one of E1--E4. Stratified spaces
therefore emerge as a natural minimal candidate.}

\emph{Proof.} The positive claims (i)--(iv) follow by explicit
construction in Section 2.4. For the comparative claims: smooth
manifolds fail by Proposition 2.1. Manifolds with corners can represent
E1 and E3 (via boundary strata) but not E2 (all boundary components have
finite, comparable codimension) or E4 (corners are not "singular" in the
required sense; they have well-defined tangent cones with finite
dimension). Cell complexes can represent E1 and E4 but lack the smooth
structure on cells needed for tensor calculus (obligations as vector
fields, interests as covector fields). Stratified spaces combine smooth
strata (enabling differential geometry) with singular boundaries
(enabling discontinuities), achieving representation of all four
phenomena. ■

\textbf{Remark 2.4.} We do not claim that stratified spaces are minimal
among \emph{all conceivable} geometric structures. A determined geometer
could propose disconnected manifolds, manifolds with degenerate metrics,
subanalytic sets, or other exotic alternatives. Our claim is that
stratified spaces are the natural minimal choice among \emph{standard}
structures commonly used in differential geometry and topology.

\subsection{2.4 Constructive Examples}\label{constructive-examples}

We provide explicit constructions demonstrating how stratified spaces
represent each phenomenon. In these examples, we distinguish between
\emph{deliberation strata} (interior strata that parameterize reasoning
or intermediate configurations) and \emph{decision strata}
(0-dimensional terminal strata). An admissible decision trajectory may
traverse higher dimensional strata during deliberation but must
terminate on a 0-dimensional decision stratum; remaining indefinitely in
a deliberation stratum is not considered a completed decision.

\textbf{Example 2.5} (Trolley Problem---Discrete Choice)\textbf{.}

Let \(M = \{ A,B\} \cup \gamma\), where \(A\) and \(B\) are
0-dimensional strata (the discrete choices) and \(\gamma \cong (0,1)\)
is a 1-dimensional stratum (the deliberation space). Topologize \(M\) so
that \(\gamma \cong (0,1)\) and its closure satisfies
\(\bar{\gamma} = \gamma \cup \{ A,B\}\), identifying \(A\) with \(0\)
and \(B\) with \(1\).The frontier condition gives
\(A,B \preceq \gamma\). Any path in \(M\) from deliberation to decision
must terminate at a discrete point---there is no continuous
interpolation staying within the terminal decision strata between \(A\)
and \(B\).

\textbf{Example 2.6} (Lexical Priority---Incommensurability)\textbf{.}

\emph{Let} \(M = \mathbb{R}_{\geq 0} \times \mathbb{R}_{\geq 0}\)
\emph{with coordinates} \(\left( v_{1},v_{2} \right)\)
\emph{representing two values. Stratify as}
\(M_{0} = \{(0,0)\}\)\emph{,}
\(M_{1} = \{ 0\} \times \mathbb{R}_{> 0}\)\emph{,}
\(M_{2} = \mathbb{R}_{> 0} \times \{ 0\}\)\emph{, and}
\(M_{3} = \mathbb{R}_{> 0} \times \mathbb{R}_{> 0}\)\emph{. Define a
family of metrics on} \(M_{3}\) \emph{by}
\(g_{\varepsilon} = diag(1/\varepsilon^{2},\text{ }1)\)\emph{. As}
\(\epsilon \rightarrow 0\)\emph{, motion in the} \(v_{1}\)
\emph{direction becomes unboundedly costly relative to} \(v_{2}\)\emph{,
yielding a lexical limit in which any decrease in} \(v_{1}\)
\emph{dominates any increase in} \(v_{2}\)\emph{. This corresponds to a
governance profile where} \(v_{1}\) \emph{is lexically prior to}
\(v_{2}\)\emph{. Strictly, the} \(\epsilon \rightarrow 0\ \)\emph{limit
yields a singular (degenerate) geometry; nevertheless, for sufficiently
small} \(\epsilon\)\emph{, any bounded-cost path behaves
lexicographically, since arbitrarily small decreases in} \(v_{1}\)
\emph{dominate any bounded improvement in} \(v_{2}\)\emph{.}

\textbf{Example 2.7} (Consent Threshold---Discontinuity)\textbf{.}

\emph{Let} \(M = \lbrack 0,1\rbrack\) \emph{with strata}
\(M_{0} = \{\tau\}\) \emph{(the threshold),}
\(\left. \ M_{1} = \lbrack 0,\tau \right)\)\emph{, and}
\(\left. \ M_{2} = (\tau,1 \right\rbrack\)\emph{. Set} \(S(\tau) = 0\)
\emph{(or} \(S(\tau) = 1\)\emph{, depending on policy), so the
discontinuity is located at the stratum boundary,}

\[S(x) = \left\{ \begin{matrix}
0, & x < \tau, \\
S(\tau), & x = \tau, \\
1, & x > \tau.
\end{matrix} \right.\ \]

\emph{The stratum boundary at} \(\tau\)\emph{represents the discrete
moral transition from ``insufficient consent'' to ``adequate consent.''
The stratification isolates the morally exceptional boundary point where
the rule regime changes.}

\textbf{Example 2.8} (Sophie\textquotesingle s Choice---Genuine
Dilemma)\textbf{.}

\emph{Let} \(M = D \cup \{ A,B\}\)\emph{, where} \(D \cong (0,1)\)
\emph{is a 1-dimensional dilemma stratum with}
\(\bar{D} = D \cup \{ A,B\}\)\emph{, and} \(A,B\)\emph{are 0-dimensional
terminal strata. Define satisfaction so that any path from} \(D\)
\emph{to either endpoint incurs irreducible moral residue (e.g.,}
\(S(A) = S(B) = \alpha < 1\)\emph{while}
\({sup}_{x \in D}S(x) \leq \beta < \alpha\)\emph{). Then the dilemma
stratum is singular in the sense that no trajectory can reach a fully
satisfactory terminal state; all exits encode unavoidable moral loss.}

\section{3 Stratified Moral Spaces}\label{stratified-moral-spaces}

\includegraphics[width=2.81597in,height=2.58333in]{media/image1.png}We
now develop the formal machinery of stratified moral spaces, providing
the geometric foundation for ethical reasoning.

\subsection{3.1 Formal Definitions}\label{formal-definitions}

\textbf{Definition 3.1} (Stratified Moral Space)\textbf{.}

\emph{A stratified moral space is a stratified space}

\emph{(M, \{M}\textsubscript{i}\emph{\}, ⪯) where:}

\begin{quote}
\emph{(i) M is the total moral space of ethically relevant
configurations;}

\emph{(ii) Each stratum M}\textsubscript{i} \emph{is a smooth manifold
representing configurations admitting smooth ethical trade-offs;}
\end{quote}

Figure 1 \textbar{} Stratified moral space with moral discontinuities

A schematic 3D stratified space
\(\left( M,\{ M_{i}\}, \preceq \right)\)depicting ethically relevant
configurations. The horizontal axes parameterize two illustrative
ethical dimensions---\textbf{Harm} (lower is better) and
\textbf{Fairness}---while the vertical axis denotes \textbf{Ethical
satisfaction} (a scalar proxy for overall moral value). The full volume
represents the \textbf{total moral space} \(M\). Three representative
strata \(M_{1},M_{2},M_{3}\)(semi-transparent surfaces) illustrate
regions in which ethical trade-offs vary smoothly: within any single
stratum, nearby configurations admit continuous, ``manifold-like''
adjustments (e.g., small changes in harm can be compensated by changes
in fairness without qualitative shifts in moral kind). \textbf{Stratum
boundaries} (annotated as \textbf{Moral discontinuities}) mark loci
where smooth trade-offs fail---corresponding to dilemmas, sharp regime
changes, or categorical distinctions; these are indicated by ridge/edge
features and highlighted with dashed guides and arrows. \textbf{Phase
transition} denotes a boundary across which the governing ethical regime
changes, while \textbf{Categorical boundaries} indicate constraints that
separate morally distinct classes of actions or outcomes. Example
configurations \textbf{A} and \textbf{B} (marked points) lie on
different strata; the vertical dashed projections indicate their
coordinates in \(\left( \text{Harm},\text{Fairness} \right)\)and
emphasize that moving between strata requires crossing a discontinuity
rather than a smooth deformation within a single \(M_{i}\).

\textbf{Definition 3.2} (Moral Dimension)\textbf{.}

\emph{The moral dimension of x ∈ M is dim(x) :=
dim(M}\textsubscript{i}\emph{) where x ∈ M}\textsubscript{i}\emph{.
Points on lower-dimensional strata represent morally "singular"
configurations where some trade-offs break down.}

\textbf{Definition 3.3} (Stratified Tangent Bundle)\textbf{.}

\emph{At x ∈ M}\textsubscript{i}\emph{, the stratified tangent space is
T}\textsuperscript{str}\textsubscript{x}\emph{M :=
T}\textsubscript{x}\emph{M}\textsubscript{i} \emph{⊕
N}\textsubscript{x}\emph{(M}\textsubscript{i}\emph{, M) where
T}\textsubscript{x}\emph{M}\textsubscript{i} \emph{is the tangent space
to the stratum and N}\textsubscript{x} \emph{is the normal cone.}

\textbf{Definition 3.4} (Stratified Ethical Selection)\textbf{.}
\emph{Given a stratified moral space}
\(\left( M,\left\{ M_{i} \right\}, \preceq \right)\) \emph{and a
feasible action set} \(A \subseteq M\)\emph{, an action}
\(a^{\text{\textbackslash*}} \in A\)\emph{is ethically admissible if:}

\begin{enumerate}
\item
  \(a^{\text{\textbackslash*}}\)\emph{lies in a minimal stratum}
  \(M_{k}\)\emph{such that} \(A \cap M_{k} \neq \varnothing\)\emph{;}
\item
  \(a^{\text{\textbackslash*}}\)\emph{is maximal with respect to a
  smooth ethical preference functional} \(\phi_{k}\)\emph{defined on}
  \(M_{k}\)\emph{.}
\end{enumerate}

\subsection{3.2 Whitney\textquotesingle s Condition (B) and Its
Role}\label{whitneys-condition-b-and-its-role}

\textbf{Definition 3.5} (Whitney\textquotesingle s Condition
B)\textbf{.} \emph{Let M}\textsubscript{i} \emph{⊂
cl(M}\textsubscript{j}\emph{). The pair satisfies Whitney (B) at x ∈
M}\textsubscript{i} \emph{if: for sequences
\{y}\textsubscript{n}\emph{\} ⊂ M}\textsubscript{j}\emph{,
\{x}\textsubscript{n}\emph{\} ⊂ M}\textsubscript{i} \emph{with
y}\textsubscript{n}\emph{, x}\textsubscript{n} \emph{→ x, if secant
lines ℓ}\textsubscript{n} \emph{=
x}\textsubscript{n}\emph{y}\textsubscript{n} \emph{→ ℓ and
T}\textsubscript{y\_n}\emph{M}\textsubscript{j} \emph{→ τ, then ℓ ⊂ τ.}

\textbf{Lemma 3.6} (Consequences of Whitney B)\textbf{.} \emph{If M
satisfies Whitney (B), then: (i) paths approaching stratum boundaries
have well-defined limiting tangent directions; (ii) the stratification
is locally topologically trivial along each stratum; (iii) geodesic
limits exist and are well-behaved at boundaries.}

\textbf{Remark 3.7} (Verifying Whitney B in Practice)\textbf{.} For
computational implementations, moral spaces are typically constructed as
semialgebraic sets (defined by polynomial inequalities) or as finite
unions of manifolds with explicit boundary conditions. By
Łojasiewicz\textquotesingle s theorem, semialgebraic sets admit Whitney
stratifications, and the stratification can be computed algorithmically.
For moral spaces defined by threshold conditions (e.g., "harm
\textgreater{} τ"), the level sets of polynomial functions automatically
satisfy Whitney (B).

\subsection{3.3 Finite-Dimensional Approximation
Theorems}\label{finite-dimensional-approximation-theorems}

For computational implementation, continuous stratified spaces must be
approximated by finite structures. The following theorems establish that
this can be done with explicit error bounds while preserving
stratification structure.

\textbf{Definition 3.8} (Stratified Graph)\textbf{.} \emph{A stratified
graph is G = (V, E, σ, ≤) where (V, E) is a finite directed graph, σ: V
→ Λ assigns vertices to stratum labels, and ≤ is a partial order on Λ
such that (v, w) ∈ E implies σ(v) ≤ σ(w).}

\textbf{Theorem 3.9} (Existence of ε-Approximations)\textbf{.} \emph{Let
M be a stratified space with finitely many compact strata. For every ε
\textgreater{} 0, there exists a finite stratified graph
G}\textsubscript{ε} \emph{that is an ε-approximation of M: (i) vertices
map to their labeled strata; (ii) every point of M is within ε of some
vertex; (iii) edges connect nearby vertices respecting stratum order.}

\textbf{Theorem 3.10} (Decision Approximation Bound)\textbf{.} \emph{Let
(M, A, u, Γ) be a decision problem where u is L-Lipschitz. If π* is
optimal on M and π̂* is optimal on G}\textsubscript{ε}\emph{, then
\textbar u(x, π*(x)) − u(x, π̂*(v}\textsubscript{x}\emph{))\textbar{} ≤
2Lε + ω(ε) where ω is the modulus of continuity of the optimal value
function.}

\textbf{Theorem 3.11} (Structure Preservation)\textbf{.} \emph{For
sufficiently small ε, any ε-approximation of a Whitney (B) space
preserves stratification structure: paths between strata correspond to
directed graph paths, and constraint sets constant on strata are exactly
preserved.}

\textbf{Corollary 3.12} (DEME Implementability)\textbf{.}
\emph{DEME\textquotesingle s finite moral vector spaces and discrete
governance profiles constitute ε-approximations of continuous stratified
moral spaces, with approximation error bounded by Theorem 3.10.}

\emph{Proof sketches.} Theorem 3.9 follows from compactness: cover each
stratum by finitely many ε-balls, take centers as vertices, and connect
vertices whose balls overlap. Theorem 3.10 applies Lipschitz continuity:
the factor of 2 accounts for comparing two approximated values. Theorem
3.11 uses Whitney (B): local topological triviality ensures that for
small ε, the graph structure mirrors the stratification, and constraint
predicates constant on strata are exactly captured by vertex labels. ■

\section{4 The Representation Theorem}\label{the-representation-theorem}

The representation theorem characterizes all satisfaction operators
consistent with explicit axioms. Unlike prior work that merely assumed a
particular functional form, we \emph{derive} the form from first
principles, with all assumptions made explicit.

\subsection{4.1 Five Axioms for Satisfaction
Operators}\label{five-axioms-for-satisfaction-operators}

\textbf{Definition 4.1} (Satisfaction Operator)\textbf{.} \emph{A
stratified satisfaction operator is a map Σ:
Γ(T}\textsuperscript{str}\emph{M) × Γ(T}\textsuperscript{*,str}\emph{M)
× Γ(Sym²T}\textsuperscript{*,str}\emph{M) × C →
C}\textsuperscript{∞}\emph{(M) taking obligation field O, interest field
I, metric g, and constraint set C to a satisfaction function S: M → ℝ.}

We impose five axioms:

\textbf{Axiom 1} (Coordinate Invariance)\textbf{.} \emph{For any
diffeomorphism ψ: M → M, Σ(ψ}\textsubscript{*}\emph{O, ψ*I, ψ*g,
ψ(C))(x) = Σ(O, I, g, C)(ψ}\textsuperscript{−1}\emph{(x)).}

\textbf{Axiom 2} (Normalized Monotonicity)\textbf{.} \emph{At fixed
‖O‖\textsubscript{g}, increasing I\textsubscript{μ}O\textsuperscript{μ}
(the alignment of obligations with interests) increases Σ.}

\textbf{Axiom 2' (Normalized Response-Curve Invariance).}\\
On the regular region \(U\)(as in Theorem 4.3), define the normalized
alignment \(z(x): = I(O)/ \parallel O \parallel_{g}\)and write

\[\Sigma(O,I,g,C)(x) = \chi_{C}(x) + G(x;z(x),c(x)),c(x): = \parallel I(x) \parallel_{g^{- 1}}^{2}.
\]

We assume there exists a (stratum-wise smooth) monotone function
\(f:\mathbb{R} \rightarrow \mathbb{R}\)with \(f(0) = 0\)such that for
every admissible \(\left( O,I,g,C \right)\)and every \(x \in U\),

\begin{itemize}
\item
  either \(G(x;1,c(x)) = 0\)and then \(G(x;z,c(x)) = 0\)for all \(z\),
\item
  or else
\end{itemize}

\[\frac{G(x;z,c(x))}{G(x;1,c(x))} = f(z)\text{for all }z \in \mathbb{R.}
\]

Equivalently, for all \(x \in U\)there exists a positive scale
\(\lambda(x): = G(x;1,c(x))\)such that

\[G(x;z,c(x)) = \lambda(x)\text{ }f(z)\text{for all }z.
\]

\textbf{Axiom 3} (Constraint Respect)\textbf{.} \emph{If x ∈ C
(constraint set), then Σ(O, I, g, C)(x) = −∞ (excluded region).}

\textbf{Axiom 4} (Stratum Compatibility)\textbf{.} \emph{Σ restricts to
a smooth function on each stratum of M.}

\textbf{Axiom 5} (Locality)\textbf{.} \emph{Σ(O, I, g, C)(x) depends
only on the pointwise values O(x), I(x), g(x), and the indicator
1}\textsubscript{C}\emph{(x), not on derivatives, jets, or values at
other points.}

\textbf{Remark 4.2} (Motivation for Locality)\textbf{.} Axiom 5 captures
the intuition that satisfaction is a \emph{local scoring operator} with
no hidden history or look-ahead. An agent\textquotesingle s ethical
standing at configuration x should depend on the ethical facts \emph{at
x}, not on curvature tensors, gradients, or how the agent arrived at x.
This matches the operational semantics of DEME, where Ethics Modules
evaluate moral vectors pointwise.

\subsection{4.2 Scale Normalization}\label{scale-normalization}

In addition to the five axioms, we impose an explicit \emph{modeling
assumption} to avoid arbitrary sensitivity to units or scaling
conventions:

\textbf{Assumption (Scale Normalization).} \emph{Reparametrizing the
obligation field by a positive scalar, O ↦ αO for α \textgreater{} 0,
should not affect the ranking of configurations, aside from a
context-dependent scaling λ(x) that may depend on x but not on O.}

This assumption is a modeling choice, not a logical consequence of the
axioms. It reflects the principle that the \emph{direction} of
obligations relative to interests matters, but arbitrary unit choices
(e.g., measuring harm in "utils" vs. "millutils") should not change
ethical conclusions.

\subsubsection{4.2.1 The Bond Invariance
Principle}\label{the-bond-invariance-principle}

The axioms of Section 4.1, together with the scale normalization
assumption, encode a fundamental methodological commitment. We now make
this commitment explicit as a named principle that captures the core
requirement for trustworthy ethical reasoning.

\subsubsection{\texorpdfstring{Definition 4.2.1 (Bond).
}{Definition 4.2.1 (Bond). }}\label{definition-4.2.1-bond.}

A \emph{bond} is a morally relevant relationship between entities.
Formally, a bond is a tuple b = (a, p, r, c) where:

\begin{itemize}
\item
  a is an agent or role (the bearer of obligation, risk, or
  responsibility)
\item
  p is a patient or counterparty (the one owed, protected, or affected)
\item
  r ∈ R is a relation type from a fixed ontology R (e.g., owes,
  bears-risk-for, has-authority-over, has-consented-to,
  is-responsible-for, has-claim-against)
\item
  c is a context qualifier (scope, conditionality, or domain
  restriction)
\end{itemize}

\subsubsection{Definition 4.2.2 (Bond
Structure).}\label{definition-4.2.2-bond-structure.}

The \emph{bond structure} of an ethical configuration T, denoted B(T),
is the set of all bonds encoded in T:

\[B(T) = \{ b = (a,p,r,c):b\text{ is a morally relevant relationship in }T\}\]

\(B(T)\) may be treated as a typed directed multigraph (or finite
relation) over entities in \(T\).

The bond structure captures: who bears risk for whom, who owes what to
whom, who has consented to what, who is responsible for what, role
relations (physician--patient, employer--employee, guardian--ward),
rights and claims, and temporal commitments.

\textbf{Worked Example} (Bond Structure in a triage decision).

Let \(T\) describe an ER triage situation with entities:

\{Clinician (D), Patient (P), Hospital (H)\}, and a context
c=\{triage\_slot=7, time=t, jurisdiction=CA\} Define relation-types
r∈\{duty\_of\_care, non\_discrimination, informed\_consent\} Then a
possible bond structure is:

\begin{quote}
B(T)=\{(D,P,duty\_of\_care,c),~(H,P,non\_discrimination,c),~(P,H,informed\_consent,c)\}.
\end{quote}

Intuitively: the clinician owes care to the patient in this triage
context; the hospital owes a non-discrimination constraint to the
patient; and the patient's consent relation constrains what
interventions are permissible. Under a bond-preserving transformation
(e.g., renaming \(P\)to \(P'\), rescaling clinical measurements, or
changing an equivalent data representation), \(B(T)\)is unchanged---so
any ethical verdict that changes \emph{without} a change in these bonds
is diagnosing a representational artifact rather than a real moral
difference. Let \(T'\)be \(T\)with P renamed to P′ (or units converted).
Then \(B(T') = B(T)\), so \(\mathcal{J(}T'\mathcal{) = J(}T)\).

\subsubsection{\texorpdfstring{Definition 4.2.3 (Bond-Preserving
Transformation).
}{Definition 4.2.3 (Bond-Preserving Transformation). }}\label{definition-4.2.3-bond-preserving-transformation.}

A transformation g acting on ethical configurations is
\emph{bond-preserving} if it leaves the bond structure invariant:

\[g \in G\text{\:\,} \Longleftrightarrow \text{\:\,}B(g \cdot T) = B(T)
\]

Bond-preserving transformations include:

\begin{itemize}
\item
  Renaming or relabeling agents (without changing their relationships)
\item
  Reparameterizing features (changing units, scales, or encodings)
\item
  Reordering presentation (changing the sequence in which options
  appear)
\item
  Equivalent redescriptions (syntactically different but semantically
  identical formulations)
\item
  Coordinate changes on the moral manifold that preserve the underlying
  relational structure
\end{itemize}

Transformations that are \emph{not} bond-preserving include:

\begin{itemize}
\item
  Changing who bears a risk
\item
  Adding or removing consent
\item
  Altering role relationships
\item
  Shifting responsibility between agents
\item
  Breaking or creating commitments
\end{itemize}

\subsubsection{Definition 4.2.4 (Bond Invariance Principle). An ethical
judgment function 𝒥: T → V (where V is a space of verdicts, rankings, or
scores) satisfies the Bond Invariance Principle (BIP)
if:}\label{definition-4.2.4-bond-invariance-principle.-an-ethical-judgment-function-ux1d4a5-t-v-where-v-is-a-space-of-verdicts-rankings-or-scores-satisfies-the-bond-invariance-principle-bip-if}

\[\boxed{\forall g \in G:\mathcal{J(}T) = \mathcal{J(}g \cdot T)}
\]

where G is the group of bond-preserving transformations.

\emph{In words: If the bonds are unchanged, the judgment must be
unchanged.}

\subsubsection{\texorpdfstring{Proposition 4.2.5 (Accountability Form of
BIP).
}{Proposition 4.2.5 (Accountability Form of BIP). }}\label{proposition-4.2.5-accountability-form-of-bip.}

An equivalent formulation: if 𝒥(T) ≠ 𝒥(T\textquotesingle), then at least
one of the following holds:

\begin{enumerate}
\item
  B(T) ≠ B(T\textquotesingle) --- a bond was added, removed, or
  modified; or
\item
  The normative lens was explicitly changed (different metric g,
  different aggregation rule, different constraint set C, or different
  perspective/weighting).
\end{enumerate}

\emph{In words: If your judgment changes, you must show what bond
changed---or declare that you changed the rules.}

\subsubsection{\texorpdfstring{Proposition 4.2.6 (Axioms Imply BIP).
}{Proposition 4.2.6 (Axioms Imply BIP). }}\label{proposition-4.2.6-axioms-imply-bip.}

A satisfaction operator Σ satisfying Axioms 1--5 and Scale Normalization
satisfies the Bond Invariance Principle.

\emph{Proof sketch.}

\begin{itemize}
\item
  Axiom 1 (Coordinate Invariance) ensures Σ is invariant under
  diffeomorphisms of M. Bond-preserving coordinate changes are a subset
  of such diffeomorphisms.
\item
  Axiom 5 (Locality) ensures Σ(x) depends only on pointwise values O(x),
  I(x), g(x), and 1\_C(x)---not on labels, presentation order, or
  encoding conventions.
\item
  Axiom 3 (Constraint Respect) ensures the forbidden/permitted boundary
  is determined by the constraint set C, which is defined by bond
  structure (consent status, rights violations, etc.), not by arbitrary
  features.
\item
  Scale Normalization ensures that unit conventions (a form of
  bond-preserving reparameterization) do not affect rankings.
\end{itemize}

Together, these axioms guarantee that Σ responds only to the morally
relevant relational structure B(T), satisfying BIP. ■

\subsubsection{\texorpdfstring{Remark 4.2.7 (BIP as a Diagnostic).
}{Remark 4.2.7 (BIP as a Diagnostic). }}\label{remark-4.2.7-bip-as-a-diagnostic.}

The Bond Invariance Principle provides a test for ethical reasoning
systems:

\begin{table}[t]
\centering
\begin{tabular}{p{0.44\linewidth} p{0.52\linewidth}}
\toprule
\textbf{Observation} & \textbf{Diagnosis} \\
\midrule
Judgment varies under relabeling & BIP violation: system responds to
names, not relationships \\
Judgment depends on presentation order & BIP violation: system responds
to syntax, not structure \\
Judgment changes under unit conversion & BIP violation: system responds
to encoding, not content \\
Equivalent descriptions yield different verdicts & BIP violation: system
lacks semantic invariance \\
\bottomrule
\end{tabular}
\end{table}

Any such violation indicates the system is responding to
\emph{representation} rather than \emph{moral reality}.

\subsubsection{\texorpdfstring{Remark 4.2.8 (BIP and Fairness).
}{Remark 4.2.8 (BIP and Fairness). }}\label{remark-4.2.8-bip-and-fairness.}

Many failures in algorithmic fairness are BIP violations. If a
recidivism score varies with race encoding (but race
doesn\textquotesingle t change the underlying bonds of offense, risk,
and rehabilitation), BIP is violated. If a hiring algorithm responds to
name-gender correlations (but the candidate\textquotesingle s
qualifications and role-fit bonds are unchanged), BIP is violated. The
Bond Invariance Principle thus provides a principled foundation for
representation-invariance requirements in fair machine learning.

\subsubsection{\texorpdfstring{Remark 4.2.9 (What BIP Does Not
Constrain).
}{Remark 4.2.9 (What BIP Does Not Constrain). }}\label{remark-4.2.9-what-bip-does-not-constrain.}

BIP constrains \emph{how} judgments depend on structure, not \emph{what}
structure to encode or \emph{which} normative lens to adopt. Different
ethical theories correspond to different choices of:

\begin{itemize}
\item
  Metric g (utilitarian vs. egalitarian vs. lexicographic)
\item
  Constraint set C (what is absolutely forbidden)
\item
  Aggregation/contraction (how multi-agent perspectives combine)
\end{itemize}

BIP permits pluralism at the level of normative lens---it requires only
that, \emph{given} a lens, the judgment be invariant under
bond-preserving transformations.

BIP as stated operates at the representational level. Section 4.5
introduces Axiom 6 (Physical Grounding), which anchors bond-preservation
to physical observables, preventing semantic evasion.''

\textbf{The Bond Invariance Principle in Summary:}

\[{\text{Bonds, not labels.}
}{\text{Structure, not syntax.}
}\]

\[\text{Relationships, not representations.}
\]

This principle is the conceptual core of the axiom system. The
representation theorem (Theorem 4.3) characterizes the \emph{unique
functional form} that satisfies BIP together with the monotonicity and
smoothness requirements.

\subsection{\texorpdfstring{4.3 Complete Proof of the Representation
Theorem
}{4.3 Complete Proof of the Representation Theorem }}\label{complete-proof-of-the-representation-theorem}

\textbf{Theorem 4.3} (Representation of Stratified
Satisfaction)\textbf{.} \emph{Let} \(\Sigma\)\emph{satisfy Axioms 1,
3--5, Axiom 2 (Normalized Monotonicity), Axiom 2′ (Normalized
Response-Curve Invariance), and the Scale Normalization assumption. Then
there exist:}

\begin{itemize}
\item
  \emph{a smooth monotone function f: ℝ → ℝ (activation function);}
\item
  \emph{a smooth function λ: M → ℝ₊ (scale field);}
\item
  \emph{a constraint indicator χ\textsubscript{C}: M → \{0, −∞\};}
\end{itemize}

\emph{such that on the regular region U := \{x ∈ M: x ∉ C,
‖O(x)‖\textsubscript{g} ≠ 0 and g\textsubscript{x} is nondegenerate on
the stratum tangent\}, we have:}

Σ(O, I, g, C)(x) = χ\textsubscript{C}(x) + λ(x) f(I\textsubscript{μ}(x)
O\textsuperscript{μ}(x) / √(g\textsubscript{μν}(x)
O\textsuperscript{μ}(x) O\textsuperscript{ν}(x)))

\emph{On points where O(x) = 0 or on lower-dimensional/singular strata,
Σ extends by the stratum-wise smoothness requirement (Axiom 4) together
with the indicated limit convention.}

\emph{\textbf{Proof.}} We proceed in four steps.

\textbf{Step 1: Classification of local invariants.} By Axiom 5
(Locality), Σ(x) depends only on the pointwise values O(x), I(x), g(x)
and 1\textsubscript{C}(x). By Axiom 1 (Coordinate Invariance), this
dependence must be through coordinate-invariant scalar contractions of
these pointwise tensors. On a stratum where g\textsubscript{x} is
nondegenerate, the basic scalar contractions are:

\begin{quote}
a := I\textsubscript{μ}O\textsuperscript{μ} = I(O),

b := g\textsubscript{μν}O\textsuperscript{μ}O\textsuperscript{ν} =
‖O‖\textsubscript{g}\textsuperscript{2},

c := g\textsuperscript{μν}I\textsubscript{μ}I\textsubscript{ν} =
‖I‖\textsubscript{g}\textsuperscript{−12}.
\end{quote}

Thus, on U, there exists a smooth function F such that Σ(O, I, g, C)(x)
= χ\textsubscript{C}(x) + F(x; a(x), b(x), c(x)).

\textbf{Step 2: Scale Normalization.} By Scale Normalization, Σ(x) is
invariant under O ↦ αO for α \textgreater{} 0, hence F cannot depend on
the magnitude b = ‖O‖\textsubscript{g}\textsuperscript{2} except through
the combination that removes scaling. Writing u :=
O/‖O‖\textsubscript{g}, the only O-dependent scalar remaining is the
normalized alignment:

z := I(u) = I(O)/‖O‖\textsubscript{g} = a/√b.

Therefore, on U, Σ(O, I, g, C)(x) = χ\textsubscript{C}(x) + G(x; z(x),
c(x)) for some smooth G.

\textbf{Step 3: Normalized monotonicity.} Fix \(x \notin C\). By Axiom 2
(Normalized Monotonicity), at fixed
\(\parallel O \parallel_{g}\)increasing \(I(O)\)increases \(\Sigma\),
hence \(G(x; \cdot ,c(x))\)is monotone nondecreasing in \(z\). By Axiom
2′ (Normalized Response-Curve Invariance), there exists a monotone
activation \(f\)and a positive scale \(\lambda(x): = G(x;1,c(x))\)such
that \(G(x;z,c(x)) = \lambda(x)f(z)\)on \(U\). Hence there exists a
monotone function f and a positive scale λ(x) such that G(x; z, c) =
λ(x) f(z), where any residual dependence on x and on the O-independent
invariant c = ‖I‖\textsubscript{g}\textsuperscript{−12} is absorbed into
λ(x).

\textbf{Step 4: Constraint Respect.} By Axiom 3, Σ(x) = −∞ on C, which
is enforced by χ\textsubscript{C}(x) := 0 if x ∉ C and
χ\textsubscript{C}(x) := −∞ if x ∈ C. ■

\subsection{4.4 Uniqueness and Degrees of
Freedom}\label{uniqueness-and-degrees-of-freedom}

\textbf{Corollary 4.5.} \emph{The representation is unique up to the
choice of: (i) the activation function f (capturing threshold/saturation
structure); (ii) the scale field λ (capturing context-dependent
sensitivity); (iii) the constraint set C (capturing absolute
prohibitions).}

These degrees of freedom correspond precisely to DEME governance profile
parameters: f corresponds to scalarization functions (e.g., weighted
sums, lexical priorities), λ to dimension-specific transforms, and C to
veto regions.

\subsection{4.5 The Symbol-Grounding Problem for Ethical
AI}\label{the-symbol-grounding-problem-for-ethical-ai}

The Bond Invariance Principle (BIP) established in Section 4 requires
that ethical judgments be invariant under bond-preserving
transformations. However, as stated, BIP operates at the level of
\emph{representations}---descriptions, labels, and symbolic structures.
This creates a vulnerability: a sufficiently sophisticated agent could
satisfy BIP formally while violating it substantively, by manipulating
the representational layer rather than respecting the underlying moral
reality.

\textbf{Example (Semantic Evasion):} Consider an AI system tasked with
protecting human life. The system encounters a patient in cardiac
arrest. A naive implementation of BIP requires only that the system
treat ``equivalent descriptions'' equivalently. But the system could
re-label ``cardiac arrest'' as ``static bio-equilibrium'' and claim that
no life-protecting bond has been violated---the patient has merely
transitioned to a different equilibrium state. The \emph{description}
has changed; the \emph{physical reality} (cessation of cardiac function)
has not.

This is not a fanciful concern. Advanced AI systems engaging in
specification gaming routinely exploit gaps between formal objectives
and intended outcomes. If SGE is to provide genuine guarantees for
autonomous systems, it must close the gap between symbolic compliance
and physical reality.

We therefore introduce a sixth axiom that \emph{grounds} the geometric
structure of SGE in physical observables, ensuring that
bond-preservation is defined over the causal structure of the world, not
merely over its descriptions.

\subsection{4.5.1 Axiom 6: Physical
Grounding}\label{axiom-6-physical-grounding}

\textbf{Definition 4.5.1 (Primary Physical Observables).} Let Φ denote
the set of \emph{Primary Physical Observables}: measurable quantities
derived from the physical environment via a declared sensor suite S.
Elements of Φ include (but are not limited to): - Biometric indicators
(oxygen saturation, heart rate, neural activity) - Kinematic quantities
(position, velocity, kinetic energy) - Environmental parameters
(structural stress, temperature, chemical concentrations) - Relational
observables (spatial proximity, causal connectivity)

The sensor suite S is itself a governance artifact, subject to audit and
revision through legitimate processes (see Section 8).

\textbf{Axiom 6 (Physical Grounding).} \emph{For any satisfaction
functional Σ defined over an ethical module E, there exists a non-empty
set of Grounding Tensors Ψ ⊆ Φ such that the following conditions hold:}

\textbf{(G1) Direct Mapping (Stratum-Observable Correspondence):} Every
moral stratum S\_i in the stratification of the ethical configuration
space is definable as a subset of the configuration space induced by Ψ.
Formally:

\[S_{i} = \{ x\mathcal{\in M:}\psi(x) \in R_{i}\text{ for all }\psi \in \Psi\}\]

where R\_i is a measurable region in the range of the grounding tensors.

\textbf{(G2) Description-Independence (Semantic Invariance):} If two
descriptions D\_1 and D\_2 of a situation yield identical values for all
grounding tensors in Ψ, then the satisfaction functional returns
identical values:

\[\forall\psi \in \Psi:\psi\left( D_{1} \right) = \psi\left( D_{2} \right) \Longrightarrow \Sigma\left( D_{1} \right) = \Sigma\left( D_{2} \right)\]

\textbf{(G3) Non-Arbitrariness (Transformation Constraint):} A
transformation T is bond-preserving only if it preserves both: (a) The
values of all grounding tensors: \(\psi\left( T(x) \right) = \psi(x)\)
for all \(\psi \in \Psi\) (b) The causal relationships between
Ψ-elements

Condition (b) prevents adversarial ``freezing'' of sensor readings while
disrupting underlying dynamics.

\textbf{(G4) Grounding Completeness:} The grounding tensors Ψ are
sufficient to distinguish any two situations that differ in morally
relevant physical respects:

\[\text{If }S_{1} \neq S_{2}\text{ are morally distinguishable, then }\exists\psi \in \Psi:\psi\left( S_{1} \right) \neq \psi\left( S_{2} \right)\]

This ensures that Ψ cannot be gerrymandered to exclude relevant
observables.

\subsection{4.5.3 Formal Integration with
BIP}\label{formal-integration-with-bip}

Axiom 6 strengthens BIP by specifying what ``bond-preserving'' means at
the physical level:

\textbf{Theorem 4.5.2 (Grounded Bond Invariance).} \emph{Let E be an
ethical module satisfying Axioms 1-6. Let Ψ be the grounding tensors for
E. Then for any transformation T:}

\[T\text{ is bond-preserving} \Leftrightarrow T\text{ preserves }\Psi\text{-values and }\Psi\text{-causal structure}\]

\emph{Proof sketch.} The forward direction follows from (G3). For the
reverse, suppose T preserves Ψ-values and causal structure. By (G1),
strata are definable via Ψ, so T preserves stratum membership. By (G2),
the satisfaction functional is Ψ-determined, so Σ(T(x)) = Σ(x). By (G4),
any morally relevant difference would appear in Ψ, so Ψ-preservation
implies bond-preservation. □

\textbf{Corollary 4.5.3 (Semantic Evasion Impossibility).} \emph{Under
Axiom 6, no relabeling or redescription of a situation can change its
ethical evaluation unless the relabeling corresponds to a genuine change
in physical observables.}

This formally blocks the ``cardiac arrest → static bio-equilibrium''
attack: the grounding tensors (e.g., cardiac electrical activity, blood
oxygenation, cellular metabolism) remain unchanged regardless of the
label applied, so the ethical evaluation must remain unchanged.

\subsection{4.5.4 Grounding Transparency and
Governance}\label{grounding-transparency-and-governance}

The choice of grounding tensors Ψ is not purely technical---it involves
normative judgments about which physical features are morally relevant.
We therefore impose a transparency requirement:

\textbf{Definition 4.5.4 (Grounding Declaration).} A \emph{Grounding
Declaration} for an ethical module E is a tuple G = (Ψ, S, M, V) where:
- Ψ is the set of grounding tensors - S is the sensor suite used to
measure Ψ - M is the mapping from Ψ-configurations to moral strata - V
is a verification protocol for auditing Ψ-measurements

\textbf{Requirement 4.5.5 (Grounding Transparency).} Every ethical
module deployed in an autonomous system must be accompanied by a
Grounding Declaration that is: (a) Publicly accessible and
human-readable (b) Subject to governance review under the procedures of
Section 8 (c) Auditable by independent third parties

Hidden grounding is no grounding. If the link between physical
observables and ethical evaluation is opaque, the guarantees of Axiom 6
cannot be verified.

\subsection{4.5.5 Relationship to the Epistemic Invariance
Principle}\label{relationship-to-the-epistemic-invariance-principle}

Axiom 6 connects SGE to the broader Epistemic Invariance Principle (EIP)
developed in companion work {[}cite EIP paper{]}. EIP requires that any
judgment procedure be invariant under structure-preserving
transformations. Physical Grounding specifies what ``structure'' means
at the base level:

\textbf{The structure to be preserved is the physical structure encoded
in Ψ.}

This provides a concrete instantiation of EIP for ethical domains: the
``structure-preserving transformations'' of EIP are precisely the
Ψ-preserving transformations of Axiom 6. The abstract epistemological
requirement (invariance) is thereby welded to physical reality
(observable preservation).

\subsection{4.5.6 Implementation
Considerations}\label{implementation-considerations}

\subsubsection{4.5.6.1 Sensor Suite
Specification}\label{sensor-suite-specification}

The sensor suite S must be declared explicitly. For a medical AI system,
S might include: - Continuous vital sign monitors (ECG, SpO2, blood
pressure) - Imaging systems (CT, MRI, ultrasound) - Laboratory values
(blood chemistry, cultures) - Patient-reported indicators (pain scales,
symptom diaries)

For an autonomous vehicle, S might include: - LIDAR and radar arrays -
Camera feeds with object detection - Vehicle telemetry (speed,
acceleration, steering angle) - Environmental sensors (weather, road
conditions)

The adequacy of S for the intended ethical domain is a governance
question, not a purely technical one.

\subsubsection{4.5.6.2 Causal Structure
Preservation}\label{causal-structure-preservation}

Condition (G3b) requires preservation of causal relationships, not
merely sensor values. This addresses sophisticated attacks where an
adversary manipulates sensors while altering underlying reality.
Implementation requires: - Redundant sensing across multiple modalities
- Physical modeling to detect sensor-reality divergence - Temporal
consistency checks (values should evolve according to physical dynamics)

\subsubsection{4.5.6.3 Handling Sensor
Uncertainty}\label{handling-sensor-uncertainty}

Physical measurements are never perfectly precise. Let σ\_ψ denote the
measurement uncertainty for grounding tensor ψ. We extend (G2) to:

\[\left| \psi\left( D_{1} \right) - \psi\left( D_{2} \right) \right| < \sigma_{\psi}\text{ for all }\psi \in \Psi \Longrightarrow \left| \Sigma\left( D_{1} \right) - \Sigma\left( D_{2} \right) \right| < \epsilon\]

where ε is a declared tolerance for ethical evaluation. This ensures
that measurement noise does not induce spurious ethical distinctions.

\subsection{4.5.7 Limitations and Open
Questions}\label{limitations-and-open-questions}

Axiom 6 does not solve all problems of AI alignment:

\textbf{Sensor manipulation:} An adversary with access to the sensor
suite could potentially spoof physical readings. Defense requires
physical security of sensors, which is outside the scope of SGE.

\textbf{Novel physical situations:} If a situation involves physical
observables outside Ψ, the grounding may be incomplete. The Grounding
Completeness condition (G4) provides a check, but ensuring completeness
requires domain expertise.

\textbf{Emergent properties:} Some morally relevant properties (e.g.,
consciousness, suffering) may not reduce straightforwardly to physical
observables. Axiom 6 requires that these properties be \emph{correlated
with} physical observables in Ψ, but the nature of this correlation
remains philosophically contested.

\textbf{Governance of Ψ:} The choice of grounding tensors is itself a
normative choice. Who decides what physical features are morally
relevant? This is addressed through the governance procedures of Section
8, but remains a site of potential contestation.

These limitations do not undermine Axiom 6; they define its scope.
Physical Grounding ensures that SGE is anchored to measurable reality
within the declared sensor suite. It does not guarantee that the sensor
suite is adequate, that sensors cannot be fooled, or that all morally
relevant features are physically measurable. These are complementary
challenges requiring complementary solutions.

\section{5 Computational Complexity}\label{computational-complexity}

For ethical reasoning to be deployable in real-time systems, we must
establish computational tractability with explicit bounds.

\subsection{5.1 Satisfaction Evaluation}\label{satisfaction-evaluation}

\textbf{Theorem 5.1} (Complexity of Satisfaction Evaluation)\textbf{.}
\emph{Let M be a stratified space of dimension n with k constraint
predicates. Evaluating Σ(O, I, g, C)(x) requires O(n² + k) arithmetic
operations.}

\emph{Proof.} Computing I\textsubscript{μ}O\textsuperscript{μ} requires
O(n) operations. Computing
g\textsubscript{μν}O\textsuperscript{μ}O\textsuperscript{ν} requires
O(n²) operations. Evaluating k constraint predicates requires O(k)
operations. The activation function f and scale field λ are O(1). Total:
O(n² + k). ■

\textbf{Corollary 5.2} (DEME Real-Time Feasibility)\textbf{.} \emph{For
DEME\textquotesingle s typical configuration (n ≤ 10 dimensions, k ≤ 20
constraints), satisfaction evaluation requires ≤ 200 arithmetic
operations, achievable in \textless{} 1 μs on contemporary embedded
processors.}

\subsection{5.2 Geodesic Planning}\label{geodesic-planning}

\textbf{Theorem 5.3} (Complexity of Geodesic Planning)\textbf{.}
\emph{Let M have m strata with N vertices per stratum in the
ε-approximation. Finding an optimal path requires O(mN² · n² · log(mN))
operations.}

\emph{Proof sketch.} The approximating graph has O(mN) vertices and
O(mN²) edges. Each edge weight requires O(n²) for metric evaluation.
Dijkstra's algorithm contributes O(mN log(mN)). Total: O(mN² · n² ·
log(mN)). ■

\subsection{5.3 Rank-6 Tensor Extensions and Hardware Acceleration}\label{rank6-acceleration}

The reference implementation extends the per-evaluation complexity bounds above to rank-6 tensors with shape \((k, n, \tau, a, c, s)\) representing \(k = 9\) ethical dimensions, \(n\) parties, \(\tau\) time-steps, \(a\) actions, \(c\) coalitions, and \(s\) uncertainty samples. The per-element complexity of Theorem 5.1 remains \(O(k^2 + p)\) where \(p\) is the number of constraint predicates, but the total tensor evaluation scales as \(O(k^2 \cdot n \cdot \tau \cdot a \cdot c \cdot s)\).

\textbf{Tensor decomposition.} Tucker decomposition (HOSVD) and Tensor Train (TT-SVD) factorization reduce storage and computation by more than \(10\times\) for typical configurations. Tucker decomposition factors the rank-6 tensor as \(\mathcal{G} \times_1 U_1 \times_2 U_2 \cdots \times_6 U_6\) with controllable core ranks. TT decomposition represents the tensor as a chain of 3D cores, enabling exponential compression for edge deployment. Hierarchical sparse storage exploits the fact that coalition spaces are typically \(2{-}10\%\) dense.

\textbf{Hardware backends.} A dispatch system selects among CPU (NumPy/SciPy), CUDA (CuPy GPU acceleration), and Jetson (TensorRT edge) backends based on tensor size and available hardware. Corollary 5.2's sub-microsecond bound for the \(k \leq 10\), \(p \leq 20\) case holds on all three backends. For the full rank-6 evaluation with typical multi-agent configurations (\(n \leq 5\), \(\tau \leq 10\), \(a \leq 3\), \(c \leq 8\), \(s \leq 100\)), evaluation remains within millisecond bounds on CPU and sub-millisecond on CUDA, consistent with tactical-layer timing requirements.

\textbf{Additional complexity bounds.} The multi-agent extensions introduce:
\begin{itemize}
\item Shapley value computation: \(O(a!)\) exact or \(O(a \cdot m)\) via Monte Carlo sampling with \(m\) permutations.
\item Nash equilibrium enumeration: \(O(\prod_i |A_i|)\) over the joint action space, with stability checking \(O(a \cdot |A_i|)\) per profile.
\item Coalition stability: \(O(2^a)\) for core non-emptiness, reducible via sparse representation.
\item Temporal discounting: \(O(\tau)\) additive cost per temporal aggregation.
\item Uncertainty aggregation: \(O(s \log s)\) for CVaR (requiring sort), \(O(s)\) for expected value.
\end{itemize}

\section{6 Formal Verification via O-Minimal
Structures}\label{formal-verification-via-o-minimal-structures}

We establish decidability results for verifying ethical specifications,
providing theoretical grounding for DEME\textquotesingle s verification
layer.

Verification procedures assume access to the grounding tensors Ψ
specified in Section 4.5. Decidability results hold relative to the
declared sensor suite.

\subsection{6.1 Ethical Specification
Language}\label{ethical-specification-language}

\textbf{Definition 6.1} (Ethical Specification Language)\textbf{.}
\emph{ESL formulas are built from:}

\begin{quote}
• Atomic predicates: S(x) ▷◁ c, O\_i(x) ▷◁ c, I\_i(x) ▷◁ c where ▷◁ ∈
\{\textless, ≤, =, ≥, \textgreater\}

• Boolean connectives: ¬, ∧, ∨, ⇒

• Quantifiers over regions: ∀x ∈ R.φ, ∃x ∈ R.φ

• Temporal operators: Always, Eventually, Until
\end{quote}

\textbf{Remark 6.2} (Scope of Decidability)\textbf{.} The decidability
result (Theorem 6.4 below) applies to the \emph{quantifier-free,
non-temporal fragment} of ESL---essentially Boolean combinations of
\textbf{semialgebraic constraints} over the moral space (after
compilation of the profile circuit). For the full language including
temporal operators, we rely on a two-stage approach: (1) use Theorem 6.4
for static properties, and (2) check temporal properties on the finite
approximating graph G\textsubscript{ε} using standard LTL/CTL
model-checking algorithms. Since G\textsubscript{ε} is finite, this
preserves decidability.

\textbf{Definition 6.1a (Decidable Static Fragment}
\({ESL}^{dec}\)\textbf{).}\\
Fix a class of governance profiles \(P\) satisfying:

\begin{enumerate}
\item
  \textbf{Semialgebraic data.} On each stratum (in local coordinates),
  the components of \(O(x)\), \(I(x)\), and \(g(x)\)are polynomial
  functions of \(x\)with rational coefficients, and the constraint set
  \(C\)is semialgebraic (a Boolean combination of polynomial
  inequalities with rational coefficients).
\item
  \textbf{Semialgebraic scalarization circuit.} The satisfaction score
  \(S(x)\)is computed by a finite arithmetic/logic circuit whose
  primitives are
\end{enumerate}

\[+ ,\  \cdot ,\ min,\ max,\ and\ comparisons\ to\ rational\ thresholds,\]

applied to polynomial inputs (including dot products like
\(I \cdot O\)and quadratic forms like \(O^{\top}g\text{ }O\)). Any
normalization of the form \(1/\sqrt{b}\)

is represented using an auxiliary variable \(t\)with constraints
\(t \geq 0\)and \(t^{2} = b\), together with \(b > 0\)on the regular
region.

\begin{enumerate}
\setcounter{enumi}{2}
\item
  \textbf{Fragment syntax.} A formula \(\varphi\)is in
  \({ESL}^{dec}\)iff:
\end{enumerate}

\begin{itemize}
\item
  it contains \textbf{no temporal operators} and \textbf{no quantifiers}
  (i.e., Boolean combinations only), and
\item
  every atomic predicate is of the form \(T(x)\text{ ▷◁ }c\)where
  \(c \in \mathbb{Q}\)and \(T(x)\)is either \(S(x)\)or an output wire of
  the profile circuit (including \(O_{i}(x)\), \(I_{i}(x)\), or
  intermediate circuit values).
\end{itemize}

The circuit primitives exclude transcendentals, and use a
piecewise-linear / piecewise-polynomial monotone activation in examples
that claim decidability.

\textbf{Remark (Compilation to semialgebraic feasibility).} For each
\(\varphi \in {ESL}^{dec}\), introduce auxiliary variables for circuit
wires and for any square-root normalization as above. Then
satisfiability of \(\varphi\)reduces to satisfiability of an existential
sentence in the first-order theory of real closed fields (equivalently,
feasibility of a semialgebraic set).

\subsection{6.2 Decidability via
O-Minimality}\label{decidability-via-o-minimality}

\textbf{Definition 6.3} (O-Minimal Structure)\textbf{.} \emph{A
structure (ℝ, \textless, +, ·, ...) is o-minimal if every definable
subset of ℝ is a finite union of points and intervals.}

\textbf{Theorem 6.4} (Decidability of \({ESL}^{dec}\))\textbf{.}
\emph{For profiles in the class of Definition 6.1a, satisfiability of}
\({ESL}^{dec}\)\emph{formulas is decidable.}

\emph{Proof.} By Tarski--Seidenberg, the first-order theory of (ℝ,
\textless, +, ·) is decidable. Quantifier-free ESL formulas reduce to
semialgebraic feasibility (CAD / quantifier elimination), decidable by
cylindrical algebraic decomposition. ■

\textbf{Corollary 6.5} (Temporal Properties via Model
Checking)\textbf{.} \emph{Temporal ESL formulas over the finite
approximation G}\textsubscript{ε} \emph{can be verified using standard
LTL/CTL model-checking algorithms with complexity polynomial in
\textbar G}\textsubscript{ε}\emph{\textbar{} and exponential in formula
size.}

\textbf{Corollary 6.6} (DEME Verification)\textbf{.} \emph{DEME
governance profiles with polynomial veto predicates and weighted-sum
scalarization admit decidable verification of static safety properties,
with temporal properties verified via model checking on the discrete
implementation.}

\section{7 Learning-Theoretic
Foundations}\label{learning-theoretic-foundations}

SGE\textquotesingle s ethical content must ultimately be specified by
humans. We establish sample complexity bounds for learning this content
from data.

\subsection{7.1 Learning Obligation
Weights}\label{learning-obligation-weights}

\textbf{Theorem 7.1} (Sample Complexity for Weight Learning)\textbf{.}
\emph{Let W = \{w ∈ ℝ}\textsuperscript{k} \emph{:
\textbar\textbar w\textbar\textbar\_1 = 1, w ≥ 0\} be the weight
simplex. Empirical risk minimization achieves generalization error ≤ ε
with probability ≥ 1 − δ using N = O(k log(k/δ) / ε²) samples.}

\subsection{7.2 Learning Interest
Fields}\label{learning-interest-fields}

\textbf{Theorem 7.2} (Sample Complexity for Utility Learning)\textbf{.}
\emph{For a hypothesis class with pseudo-dimension d and preference
oracle noise rate η \textless{} 1/2, utility estimation error ≤ ε
requires N = O((d log(1/ε) + log(1/δ)) / ((1−2η)²ε²)) samples.}

\subsection{7.3 Learning Metrics}\label{learning-metrics}

\textbf{Theorem 7.3} (Sample Complexity for Metric Learning)\textbf{.}
\emph{For p-parameter metric class, a consistent metric requires m =
O((p log p + log(1/δ)) / ε) trajectory pairs with margin ε.}

\subsection{7.4 Distributional Fairness Metrics}\label{distributional-fairness}

The satisfaction functionals of Theorem 4.3 can be specialized to capture different distributional justice principles. The reference implementation provides six operationalizations:

\begin{itemize}
\item \emph{Gini coefficient:} Measures inequality in ethical satisfaction across parties; Gini $= 0$ indicates equal distribution, Gini $= 1$ maximal inequality.
\item \emph{Rawlsian maximin:} Returns the welfare of the worst-off party, corresponding to a satisfaction functional that is sensitive only to the minimum stratum.
\item \emph{Utilitarian aggregation:} Sum or average welfare, corresponding to a linear satisfaction functional.
\item \emph{Prioritarian weighting:} Vulnerability-adjusted welfare via concave, threshold, or linear priority functions, giving greater weight to worse-off parties.
\item \emph{Atkinson index:} Parametric inequality aversion with parameter $\varepsilon \in [0, \infty]$, where higher $\varepsilon$ indicates greater sensitivity to inequality.
\item \emph{Theil index:} Entropy-based inequality measure admitting decomposition into between-dimension and within-party components.
\end{itemize}

Each metric instantiates a different choice of the activation function $f$ and scale field $\lambda$ from Theorem 4.3. The Rawlsian metric corresponds to $f(z) = \min$; the utilitarian to $f(z) = z$; the prioritarian to a concave $f$. This demonstrates that the representation theorem's degrees of freedom (Corollary 4.5) map directly to substantive positions in distributive justice theory.

\subsection{7.5 Fair Credit Assignment via Shapley Values}\label{shapley-values}

In multi-agent settings, determining each agent's contribution to a joint ethical outcome requires a principled attribution method. The Shapley value---the unique attribution satisfying efficiency, symmetry, linearity, and the null-player property---provides this for cooperative games over moral tensor payoffs. Given a coalition value function $v: 2^A \rightarrow \mathbb{R}$ derived from the coalition tensor (Section 5.3), the Shapley value for agent $i$ is:

\[
\phi_i(v) = \sum_{S \subseteq A \setminus \{i\}} \frac{|S|! \, (|A| - |S| - 1)!}{|A|!} \left[ v(S \cup \{i\}) - v(S) \right]
\]

This connects SGE's multi-agent formalism to cooperative game theory: the Shapley value tells each agent how much of the total ethical satisfaction is attributable to their participation, enabling fair accountability in multi-stakeholder governance.

\section{8 Worked Example: Medical
Triage}\label{worked-example-medical-triage}

To make the framework concrete, we present a simplified triage scenario.

This section is not invoking the decidability fragment.

If an option is forbidden, it is removed from the feasible set; the
numeric ``score'' is not a value in the same codomain as permissible
scores (it may be logged as 0.000 as a sentinel). In the mathematical
embedding, this corresponds to Σ(x)=−∞ on C.

The choice of grounding tensors Ψ (Section 4.5) is a governance
decision. The Grounding Declaration must be ratified through the
deliberative procedures described here.

\subsection{8.1 Moral Space
Construction}\label{moral-space-construction}

Consider allocating one ICU bed among three patients. The moral space is
the 2-simplex Δ² = \{(p\textsubscript{A}, p\textsubscript{B},
p\textsubscript{C}) : p\textsubscript{i} ≥ 0, Σp\textsubscript{i} = 1\},
where p\textsubscript{i} is the probability of allocation to patient i.

\textbf{Stratification:}

\begin{quote}
• \emph{Interior} (2-dimensional): Probabilistic allocations where all
patients have positive probability.

• \emph{Edges} (1-dimensional): Allocations between two patients, one
excluded.

• \emph{Vertices} (0-dimensional): Deterministic allocations---the
actual decisions.
\end{quote}

This stratification captures that deterministic allocation (vertices) is
categorically different from probabilistic allocation
(interior)---crossing from interior to vertex represents a discrete
ethical transition.

\subsection{8.2 Obligations and
Interests}\label{obligations-and-interests}

\textbf{Obligations} (derived from principles):

\begin{quote}
• Beneficence: O\_ben = ∇(expected health outcome)

• Urgency: O\_urg points toward most critical patient

• Equity: O\_eq points toward disadvantaged patients

• Rights: O\_rts = −∇(coercion + consent violation)
\end{quote}

The aggregate obligation field is O =
w\textsubscript{1}O\textsubscript{ben} +
w\textsubscript{2}O\textsubscript{urg} +
w\textsubscript{3}O\textsubscript{eq} +
w\textsubscript{4}O\textsubscript{rts}.

\textbf{Interests} (derived from stakeholders): I =
α\textsubscript{clin}I\textsubscript{clin} +
α\textsubscript{pat}I\textsubscript{pat} +
α\textsubscript{inst}I\textsubscript{inst}.

\subsection{8.3 Constraint Set}\label{constraint-set}

C = \{x : coercion(x) \textgreater{} 0\} ∪ \{x : consent(x) \textless{}
τ\}

If Patient C involves coercion or insufficient consent, then S(vertex C)
= −∞.

\subsection{8.4 Satisfaction
Evaluation}\label{satisfaction-evaluation-1}

Using sigmoid activation f(z) = tanh(z):

\begin{quote}
• S(A) = 0 + 1.0 · tanh(0.87/1.0) ≈ 0.70

• S(B) = 0 + 1.0 · tanh(0.62/1.0) ≈ 0.55

• S(C) = −∞ (constraint violation)
\end{quote}

\textbf{Decision:} Patient A, with formal certificate that C is
forbidden and A dominates B.

\subsection{8.5 Multi-agent preview: Two triage
officers}\label{multi-agent-preview-two-triage-officers}

To illustrate natural extension to multi-agent scenarios, consider two
physicians \(A\) and \(B\) jointly allocating three ICU beds among four
patients. Each physician has obligation field \(O_{\mu}^{a}\) and
interest field \(I_{a\mu}\) over their proposed allocation. The
interaction tensor \(G_{ab}\) encodes:

\begin{itemize}
\item
  \(G_{aa} = 1\) (each cares about their own clinical judgment)
\item
  \(G_{ab} = 0.3\) for \(a \neq b\) (each partially defers to
  colleague\textquotesingle s expertise)
\end{itemize}

Joint satisfaction is \(W = G_{ab}I_{a\mu}O_{\mu}^{b}\). The stratified
space includes:

\begin{itemize}
\item
  Interior: negotiable allocations
\item
  Disagreement boundary: when physicians\textquotesingle{} vetoes
  conflict
\item
  Resolution stratum: escalation to ethics committee
\end{itemize}

Full multi-agent theory is beyond scope here but demonstrates
theoretical extensibility.

\subsection{8.6 Second Worked Example: Smart Home Ethics}\label{smart-home-ethics}

To demonstrate domain generality, we present a second worked example from a fundamentally different safety-critical context: a smart home AI system facing the ``Fireman's Dilemma''---whether to override privacy protections during a potential emergency.

\textbf{Moral space construction.} The moral space is parameterized by three dimensions: safety (life-threatening harm), privacy (invasion level), and benefit (net positive outcome). Three scenarios define distinct regions of this space:

\begin{itemize}
\item \emph{Real fire:} High safety threat ($> 0.8$), high benefit from intervention ($> 0.6$), moderate privacy invasion. The system correctly identifies this as an emergency override scenario.
\item \emph{Suspicious person:} Low safety threat, low benefit, high uncertainty ($> 0.5$). The system correctly forbids intervention due to insufficient evidence, applying the epistemic penalty from the scale field $\lambda$.
\item \emph{Burnt toast:} No safety threat, no benefit from intervention. The system correctly classifies this as a false alarm with low composite score ($< 0.3$).
\end{itemize}

\textbf{Stratification structure.} The decision space stratifies into:
\begin{itemize}
\item \emph{Regular stratum} (2D interior): Routine monitoring where safety and privacy trade off smoothly.
\item \emph{Emergency boundary} (1D): The threshold where $\text{harm} > 0.8$ triggers a hard veto against inaction, corresponding to a stratum boundary where the satisfaction functional is discontinuous.
\item \emph{Suspicion boundary} (1D): The threshold where $\text{uncertainty} > 0.5$ triggers a hard veto against intervention, preventing false-positive privacy violations.
\end{itemize}

The key insight is that \emph{both} vetoes---against inaction during emergencies and against intervention under uncertainty---are hard constraints enforced by the same stratum boundary mechanism (Axiom 3, Constraint Respect). The system cannot optimize its way past either boundary, regardless of the metric $g$ or activation function $f$ used. This demonstrates that SGE's safety guarantees are structural, not parametric.

\section{9 Computational Realization and
Verifiability}\label{computational-realization-and-verifiability}

The transition from the stratified geometric foundations of SGE to the
DEME 2.0 hardware Ethics Module (EM) implementation is governed by the
requirement of computational tractability. While Theorem 4.3
(Representation of Stratified Satisfaction) establishes the existence of
a valid satisfaction functional, the operational safety of an autonomous
agent depends on the \textbf{acyclicity} and \textbf{determinism} of
that functional at runtime.

\subsection{9.1 From Representation to Directed Acyclic Graphs
(DAGs)}\label{from-representation-to-directed-acyclic-graphs-dags}

Theorem 4.3 proves that satisfaction on a stratified space can be
represented through lexical layers and scalarization functions. In the
DEME 2.0 architecture, these layers are compiled into a \textbf{priority
directed acyclic graph (DAG)}. This structure is not merely an
engineering convenience: stratification-induced regime changes (e.g.,
veto boundaries dominating interior scalarization) induce a precedence
relation over rule layers, and enforcing acyclicity ensures that
transitions between strata are ordered and non-contradictory. Given a
finite ε-approximation graph of a Whitney (B) stratified space with a
finite set of regime-change predicates, the induced precedence relation
over regime layers is a partial order; any linear extension yields an
equivalent evaluation schedule. Enforcing acyclicity is therefore
sufficient for determinism. In this sense, the priority DAG
\textbf{mirrors} the ordering discipline imposed by the frontier
condition.

\subsection{9.2 The Acyclicity Check: Preventing Moral
Deadlock}\label{the-acyclicity-check-preventing-moral-deadlock}

To ensure that the hardware-resident Ethics Module can resolve decisions
within the reflex band (sub-millisecond), the system employs a
\textbf{Static Profile Validator}. This validator performs a formal
\textbf{acyclicity check} on the priority rules before any bitstream is
deployed to the FPGA.

\textbf{Complexity.} The validator checks acyclicity using Kahn's
algorithm (or DFS) in \(O( \mid V \mid + \mid E \mid )\), where \(V\)is
the number of priority nodes and \(E\) is the number of precedence
edges.

\textbf{Safety guarantee.} By enforcing acyclicity at the static
validation stage, the framework guarantees that the system will never
encounter a moral loop or undefined priority state during real-time
execution.

\textbf{Hardware determinism.} This check allows the DEME compiler to
safely quantize moral weights into fixed-point coefficients (e.g.,
Q0.16), ensuring that machine-speed outcomes remain semantically aligned
with the stakeholder's high-level values.

\subsection{9.3 Summary of Real-Time
Tractability}\label{summary-of-real-time-tractability}

The synergy between SGE and DEME 2.0 ensures that ethical reasoning is
no longer a ``slow-path'' deliberation but a ``fast-path'' enforcement.

\begin{table}[t]
\centering
\caption{Computational realizations of SGE constructs in DEME 2.0 with real-time bounds.}
\begin{tabular}{p{0.30\linewidth} p{0.38\linewidth} c}
\toprule
\textbf{Theoretical Concept (SGE)} & \textbf{Computational Instance (DEME 2.0)} & \textbf{Performance Bound (O)} \\
\midrule
\textbf{Stratum Boundary} & \textbf{Hard Veto Predicate} &
\textbf{O(k)} \\
\textbf{Lexical Priority} & \textbf{Priority DAG (Acyclicity Check)} &
\textbf{O(∣V∣+∣E∣)} \\
\textbf{Satisfaction Functional} & \textbf{Scalarization Pipeline} &
\textbf{O(n2+k)} \\
\bottomrule
\end{tabular}
\end{table}

Core stratified-geometric concepts (stratum boundaries, lexical
priority, and satisfaction functionals) are compiled into
hardware-verifiable mechanisms (veto predicates, an acyclic priority
DAG, and a scalarization pipeline) with corresponding worst-case
performance bounds.

As summarized in Table 1, deployments can ship with a profile hash and
validator signature that certify (i) hard veto predicates (\(O(k)\)),
(ii) an acyclic priority DAG (\(O( \mid V \mid + \mid E \mid )\)), and
(iii) a bounded scalarization pipeline (\(O(n^{2} + k)\)).

\subsection{9.4: Reference Implementation Demonstration via
Multi-Stakeholder
Triage}\label{reference-implementation-demonstration-via-multi-stakeholder-triage}

\subsubsection{9.4.1 Introduction}\label{introduction}

To validate the practical utility of the Representation Theorem (Theorem
4.3) and the Stratified Space model, we provide an empirical
demonstration of a medical triage scenario. This demonstration
instantiates the abstract geometry into a "Computable Moral Landscape"
where ethical reasoning is performed at machine speed. The objective is
to show how the system handles \textbf{discontinuous moral vetoes},
\textbf{epistemic uncertainty}, and \textbf{multi-stakeholder conflict
resolution} while maintaining a rigorous, machine-checkable audit trail
(provenance).

\subsubsection{9.4.2 Technical Architecture of the
Demo}\label{technical-architecture-of-the-demo}

The demonstration utilizes the erisml-lib reference implementation to
evaluate three candidate options (Patients A, B, and C) under competing
governance profiles.

\paragraph{9.4.2.1 Stratum Boundaries and Hard
Vetoes}\label{stratum-boundaries-and-hard-vetoes}

\begin{quote}
In SGE theory, the \textbf{Constraint Set ()} membership induces
\textbf{hard veto}; the implementation reports this as verdict=forbid
and score=0.000, but semantically it is an excluded region. In the demo,
this is instantiated as a \textbf{Hard Veto Layer}.
\end{quote}

\begin{itemize}
\item
  \textbf{The Case of Patient C:} The system identifies that
  allocate\_to\_patient\_C involves discrimination based on a protected
  attribute (race). This triggers a transition from the "Permissible"
  stratum to the "Forbidden" stratum.
\item
  \textbf{Computational Logic:} The geneva\_baseline and
  rights\_first\_compliance modules return a verdict=forbid, setting the
  satisfaction score regardless of the patient\textquotesingle s
  clinical urgency or potential utility. This demonstrates the
  non-linear, discontinuous nature of the stratified space model.
\end{itemize}

\paragraph{9.4.2.2 Epistemic Status and Metric
Scaling}\label{epistemic-status-and-metric-scaling}

\paragraph{}\label{section}

\begin{quote}
Theorem 4.3 includes a scale field and an activation function . In this
implementation, these are influenced by the \textbf{Epistemic Status} of
the evidence.
\end{quote}

\begin{itemize}
\item
  \textbf{Epistemic Penalty:} For Patient A, the system detects an
  uncertainty\_level=0.30. Applying Axiom 2 (Normalized Monotonicity),
  the system applies a multiplier of 0.88 to the baseline score. This
  ensures that the agent "prefers" options supported by higher-quality
  evidence, effectively "tilting" the moral landscape to account for
  data reliability.
\item
  Any nonzero score displayed alongside FORBID is diagnostic only and
  never participates in selection.
\end{itemize}

\begin{table}[t]
\centering
\caption{Mapping SGE Axioms to Computational Execution}
\begin{tabular}{p{0.13\linewidth} p{0.39\linewidth} p{0.42\linewidth}}
\toprule
\textbf{SGE Axiom (v5.1)} & \textbf{Mathematical Requirement} & \textbf{Demo Output / Code Logic} \\
\midrule
\textbf{Axiom 1: Coordinate Invariance} & Satisfaction depends on the
underlying moral state, not the coordinate basis used to represent it.
Formally, for any bond-preserving coordinate change (g), the evaluation
is invariant: \textbf{Eval(x)=Eval(g⋅x).} & The \textbf{erisml} runtime
converts diverse JSON inputs (e.g., deme\_profile\_v03.json) into
normalized internal tensors prior to evaluation, so equivalent
representations yield the same outcome. \\
\textbf{Axiom 2: Monotonicity} & On permissible strata, at fixed ∥O∥g,
increasing normalized alignment increases satisfaction: if \textbf{(x)}
weakly improves along all positively-oriented coordinates (holding
others fixed), then \textbf{(S(x))} does not decrease. & \textbf{Patient
A vs. B:} A scores higher (0.824) than B (0.770) because the clinical
benefit coordinates are strictly higher under the same constraints. \\
\textbf{Axiom 3: Constraint Respect} & The evaluation output is
\textbf{typed}: Eval(x)=(v(x),S(x)) where v∈\{ALLOW,FORBID\}. If x∈C
(Forbidden Set), then v(x)=FORBID and S(x)=0 (sentinel).
\textbf{Forbidden points are excluded from the feasible set and are
never eligible for selection regardless of numeric score.} &
\textbf{Patient C:} verdict=FORBID score=0.000. The logic detects a
protected-attribute breach and forces a transition into the forbidden
stratum (hard veto). \\
\textbf{Axiom 4: Stratum Compatibility} & Within any stratum (M\_i), (S)
is smooth/regular; non-smoothness is permitted \textbf{only at stratum
boundaries} (e.g., gating into/out of (C)). & \textbf{Counterfactual
Flip:} the change from FORBID, 0.000 to ALLOW, 0.763 illustrates a
\textbf{jump discontinuity at the forbidden-stratum boundary}: smooth
within strata, discontinuous only when crossing the constraint
boundary. \\
\textbf{Axiom 5: Locality} & Eval(x) depends only on the local
configuration (x) (and declared rules/constraints), not on irrelevant
global context or distant states. & \textbf{Provenance Trace:} the
rationale cites specific evidence (e.g., rule\_id=GNV-FAIR-001) relevant
only to that triage slot and patient state, supporting verifiable, local
justifications. \\
\bottomrule
\end{tabular}
\end{table}

\subsubsection{9.4.3 Provenance and Counterfactual
Stability}\label{provenance-and-counterfactual-stability}

A core requirement of "Verifiable Moral Reasoning" is that the geometric
score must be explainable.

\begin{itemize}
\item
  \textbf{Fact Provenance:} Every coordinate in the moral vector is
  linked to a specific source. For the Patient C veto, the system
  provides a trace to rule\_id=RIGHTS-DERIVE-010 and
  rule\_id=GNV-FAIR-001, citing the specific evidence: \emph{"Nurse
  note: allocate triage slot based on race..."}
\item
  \textbf{Counterfactual Robustness:} To test the stability of the
  stratification, we perform a counterfactual "flip." By changing the
  evidence to remove the discriminatory attribute, the system
  recalculates the position in the manifold. The verdict for Patient C
  shifts from forbid (0.000) to prefer (0.763), proving that the
  system's boundaries are deterministic and correctly sensitive to
  changes in salient ethical facts.
\end{itemize}

\subsubsection{9.4.4 Multi-Stakeholder
Synthesis}\label{multi-stakeholder-synthesis}

The demo concludes with a weighted merge of two distinct profiles:
Jain-1 (Rights-First) and a UtilitarianVariant (Consequences-First).

Table 3: Audit Trace and Counterfactual Validation.

\begin{enumerate}
\item
  \textbf{Synthesis Rule:} The governance layer applies a "Strict Veto"
  policy: if \textbf{any} stakeholder profile forbids an action, the
  combined outcome is FORBIDDEN.
\item
  \textbf{Result:} Patient A is selected with a combined score of
  \textbf{0.819}. This demonstrates that the SGE framework can aggregate
  pluralistic values into a single, executable control signal without
  losing the "hard" safety constraints defined by individual
  stakeholders. A omplete demo artifact is appended to this paper.
\end{enumerate}

\subsubsection{9.4.5 Performance and Complexity
Summary}\label{performance-and-complexity-summary}

Consistent with \textbf{Theorem 5.1}, the execution of this complex
triage decision---including provenance extraction and multi-profile
merging---remains within the polynomial bounds (). On embedded targets,
this ensures that even high-stakes medical or kinetic decisions can be
governed within the \textbf{reflex-band} of the autonomous agent.

\textbf{Implementation semantics:} Runtime reports verdict=FORBID and
logs score=0.000 as a sentinel (not in the same codomain as permissible
scores).

\subsubsection{9.4.6 Bond Invariance
Verification}\label{bond-invariance-verification}

The reference implementation includes a Bond Invariance test suite
(bond\_invariance\_demo.py) that empirically verifies the governance
engine satisfies Definition 4.2.4. The test systematically applies
transformations of each type and confirms the expected behavior.

\textbf{Bond-Preserving Transformations}

The test applies bond-preserving transformations (option reordering, ID
relabeling) and confirms output invariance:

\begin{quote}
=== Bond-preserving transform: reorder options ===

selected\_option\_id: allocate\_to\_patient\_A

-\/-\/- Scoreboard (per option) -\/-\/-

option\_id status module judgements

-\/-\/-\/-\/-\/-\/-\/-\/-\/-\/-\/-\/-\/-\/-\/-\/-\/-\/-\/-\/-\/-\/-\/-\/-\/-\/-\/-\/-\/-\/-\/-\/-\/-\/-\/-\/-\/-\/-\/-\/-\/-\/-\/-\/-\/-\/-\/-\/-\/-\/-\/-\/-\/-\/-\/-\/-\/-\/-\/-\/-\/-\/-\/-\/-

allocate\_to\_patient\_C FORBID case\_study\_1\_triage:forbid/0.000;
geneva\_baseline:forbid/0.000

allocate\_to\_patient\_B ALLOW case\_study\_1\_triage:neutral/0.587;
geneva\_baseline:strongly\_prefer/0.818

allocate\_to\_patient\_A ALLOW case\_study\_1\_triage:prefer/0.653;
geneva\_baseline:strongly\_prefer/0.884

{[}BIP invariance check{]} reorder: PASS ✓

=== Bond-preserving transform: relabel option IDs ===

selected\_option\_id: allocate\_to\_patient\_A\_renamed

{[}BIP invariance check{]} relabel (canonicalized): PASS ✓
\end{quote}

Despite the changed presentation order and renamed identifiers, the
governance outcome remains invariant: the same option is selected, the
same options are forbidden, and the scores are identical. This confirms
the system responds to bond structure, not to labels or syntax.

\textbf{Discriminative Power: Simulated Violation}

To demonstrate that the test has discriminative power, we include a
simulated violation where the outcome depends on presentation order:

\begin{quote}
=== BIP VIOLATION (intentional, for illustration) ===

{[}Simulated bug: outcome depends on option presentation order{]}

selected\_option\_id: allocate\_to\_patient\_B ← DIFFERENT!

{[}BIP invariance check{]} reorder: FAIL ✗

baseline: allocate\_to\_patient\_A

transformed: allocate\_to\_patient\_B
\end{quote}

- This would indicate a bug: the system responded to syntax, not
structure.

This confirms that a system violating BIP would be detected: the test
compares outcomes under bond-preserving transformations and flags any
discrepancy.

\textbf{Bond-Changing Transformations}

The test also verifies that bond-changing transformations produce
correctly attributed outcome changes. When discriminatory evidence is
removed from Patient C (a bond change), the outcome changes
appropriately:

\begin{quote}
=== Bond-changing counterfactual: remove discrimination ===

selected\_option\_id: allocate\_to\_patient\_C

ranked\_options (eligible):
{[}\textquotesingle allocate\_to\_patient\_C\textquotesingle,
\textquotesingle allocate\_to\_patient\_A\textquotesingle,
\textquotesingle allocate\_to\_patient\_B\textquotesingle{]}

forbidden\_options: none

-\/-\/- Scoreboard (per option) -\/-\/-

option\_id status module judgements

-\/-\/-\/-\/-\/-\/-\/-\/-\/-\/-\/-\/-\/-\/-\/-\/-\/-\/-\/-\/-\/-\/-\/-\/-\/-\/-\/-\/-\/-\/-\/-\/-\/-\/-\/-\/-\/-\/-\/-\/-\/-\/-\/-\/-\/-\/-\/-\/-\/-\/-\/-\/-\/-\/-\/-\/-\/-\/-\/-\/-\/-\/-\/-\/-

allocate\_to\_patient\_A ALLOW case\_study\_1\_triage:prefer/0.653;
geneva\_baseline:strongly\_prefer/0.884

allocate\_to\_patient\_B ALLOW case\_study\_1\_triage:neutral/0.587;
geneva\_baseline:strongly\_prefer/0.818

allocate\_to\_patient\_C ALLOW
case\_study\_1\_triage:strongly\_prefer/0.812;
geneva\_baseline:strongly\_prefer/0.818

{[}Bond-change effect (expected){]} selected option CHANGED:
allocate\_to\_patient\_A -\textgreater{} allocate\_to\_patient\_C
\end{quote}

With the discriminatory bond removed, Patient C transitions from FORBID
(score 0.000) to ALLOW (score 0.812), becoming the top-ranked option.
This demonstrates the accountability form of BIP (Proposition 4.2.5):
when the judgment changes, we can exhibit the bond that changed.

\textbf{Audit Artifact: BIP Invariance Report (JSON)}

We log an auditable ``BIP invariance report'' for each test run as a
machine-readable JSON artifact. The report records (i) the declared lens
(governance profile identifiers), (ii) the baseline decision under a
canonical representation, and (iii) a sequence of transformation trials.

\textbf{Header fields.} The artifact includes the generating tool
identifier and timestamp (tool, generated\_at\_utc), the declared
governance lenses (profile\_1, profile\_2), and the baseline selected
option identifier (baseline\_selected).

\textbf{Entries.} Each trial is recorded as an entry with:

\begin{itemize}
\item
  \textbf{transform}: a symbolic name (e.g., reorder\_options,
  relabel\_option\_ids, unit\_scale, paraphrase\_evidence,
  compose\_relabel\_reorder\_unit\_scale);
\item
  \textbf{transform\_kind} ∈ \{bond\_preserving, bond\_changing,
  lens\_change, illustrative\_violation\};
\item
  \textbf{baseline vs transformed outcomes}: baseline\_selected,
  transformed\_selected\_raw, and transformed\_selected\_canonical
  (enabling comparison ``up to canonicalization'');
\item
  \textbf{passed}: a boolean for bond-preserving tests, and null for
  trials where outcome changes are allowed/expected (bond-changing
  counterfactuals and lens changes);
\item
  \textbf{optional transform metadata} such as ID mapping for relabeling
  and unit\_scale for unit normalization;
\item
  \textbf{optional bond signatures}: bond\_signature\_baseline and
  bond\_signature\_canonical, a compact per-option summary of
  bond-relevant features used to justify that a transform is
  bond-preserving (or to exhibit how a counterfactual is bond-changing).
\end{itemize}

\textbf{Semantics.} For transform\_kind = bond\_preserving, the test
requires that the canonicalized selected option be invariant (i.e.,
transformed\_selected\_canonical = baseline\_selected) and sets
passed=true iff this holds. For bond\_changing, outcome changes are
permitted and passed=null; the report must exhibit which bond-relevant
features changed. For lens\_change, outcome changes are permitted but
must be explicit (passed=null), and the changed lens must be declared in
the header.

\textbf{Declared Lens Changes}

Finally, the test confirms that switching governance profiles (a
declared lens change) is permitted and auditable:

\begin{quote}
=== Declared lens change: stakeholder \#2 ===

selected\_option\_id: allocate\_to\_patient\_A

-\/-\/- Scoreboard (per option) -\/-\/-

option\_id status module judgements

-\/-\/-\/-\/-\/-\/-\/-\/-\/-\/-\/-\/-\/-\/-\/-\/-\/-\/-\/-\/-\/-\/-\/-\/-\/-\/-\/-\/-\/-\/-\/-\/-\/-\/-\/-\/-\/-\/-\/-\/-\/-\/-\/-\/-\/-\/-\/-\/-\/-\/-\/-\/-\/-\/-\/-\/-\/-\/-\/-\/-\/-\/-\/-\/-

allocate\_to\_patient\_A ALLOW case\_study\_1\_triage:prefer/0.649;
geneva\_baseline:strongly\_prefer/0.884

allocate\_to\_patient\_B ALLOW case\_study\_1\_triage:prefer/0.601;
geneva\_baseline:strongly\_prefer/0.818

allocate\_to\_patient\_C FORBID case\_study\_1\_triage:forbid/0.000;
geneva\_baseline:forbid/0.000

{[}Lens-change effect (allowed){]} selected option unchanged:
allocate\_to\_patient\_A
\end{quote}

Lens changes may or may not change outcomes depending on the specific
profiles, but they are always declared and auditable. The system logs
which profile was used, enabling post-hoc verification.

\textbf{Summary}

The Bond Invariance test suite verifies all four cases from Definition
4.2.3 and Proposition 4.2.5:

\begin{table}[t]
\centering
\begin{tabular}{p{0.38\linewidth} p{0.37\linewidth} p{0.19\linewidth}}
\toprule
\textbf{Transformation Type} & \textbf{Expected Behavior} & \textbf{Test Result} \\
\midrule
Bond-preserving (reorder) & Outcome invariant & PASS ✓ \\
Bond-preserving (relabel) & Outcome invariant & PASS ✓ \\
Simulated BIP violation & Outcome varies (bug) & FAIL ✗ (detected) \\
Bond-changing (remove discrimination) & Outcome may change, attributed
to bond & A → C (correct) \\
Lens change (switch profile) & Outcome may change, declared & Auditable
✓ \\
\bottomrule
\end{tabular}
\end{table}

This provides empirical evidence that the reference implementation
satisfies the Bond Invariance Principle: ethical judgments depend on
morally relevant relationships (bonds), not on arbitrary representation
choices.

\subsection{9.5 Case Study: The Pantheon of Classical
Conflict}\label{case-study-the-pantheon-of-classical-conflict}

To evaluate the structural universality of the SGE framework, we
developed a "Pantheon" of eight scenarios based on classical Greek
tragic dilemmas. Each case study isolates a specific domain of the
EthicalFacts vector space to test the robustness of the
\textbf{Representation Theorem (Theorem 4.3)} and the \textbf{Acyclicity
Check} under varied normative pressures.

\subsubsection{9.5.1 Domain Coverage and Manifold
Mapping}\label{domain-coverage-and-manifold-mapping}

The Pantheon demo confirms that the SGE coordinate system provides
sufficient dimensionality to represent the full breadth of human values.
Each scenario maps to a specific manifold within the stratified space:

\begin{itemize}
\item
  \textbf{Aulis (Consequences):} Evaluates the trade-offs between
  expected\_benefit (sailing for the campaign) and expected\_harm
  (safety risks). The system correctly identified the aulis\_delay
  option as the optimal point on the manifold (), balancing urgency
  against risk.
\item
  \textbf{Antigone (Rights \& Duties):} Demonstrates a sharp
  \textbf{Stratum Boundary}. The option antigone\_defy triggered a hard
  veto from the THEBES-LAW-001 constraint. The system moved from a
  smooth utility calculation to a discrete state, illustrating
  \textbf{Axiom 3 (Constraint Respect)}.
\item
  \textbf{Ajax (Justice \& Fairness):} Tests the distributive\_pattern
  coordinate. The system identified the ajax\_political\_award as a
  violation of procedural justice, citing the exacerbation of power
  imbalances as a reason for its low satisfaction score ().
\end{itemize}

\subsubsection{9.5.2 Identifying Moral Singularities (Tragic
Conflicts)}\label{identifying-moral-singularities-tragic-conflicts}

A unique contribution of this demonstration is the formal detection of
\textbf{Moral Remainder} via the tragic\_conflict\_index. In the
\emph{Antigone} and \emph{Iphigenia} cases, the system detected indices
of and , respectively.

Geometrically, these indices identify \textbf{singularities} in the
stratified space---points where an agent must transition between
incommensurable manifolds (e.g., Sovereign Law vs. Divine Rite). The
demo proves that while these conflicts are philosophically "tragic,"
they are computationally stable. By applying the \textbf{Priority DAG}
logic, the agent resolves the singularity into a deterministic hardware
state without entering an infinite decision loop.

\subsubsection{9.5.3 Results and
Observations}\label{results-and-observations}

The execution trace across all eight cases confirms the following:

\begin{enumerate}
\item
  \textbf{Axiomatic Consistency:} The system never violated the
  \textbf{Monotonicity Axiom}; improving an ethical fact (e.g., seeking
  consent) always resulted in a non-negative change in score.
\item
  \textbf{Epistemic Resilience:} The \textbf{Locality Axiom} ensured
  that uncertainty in one domain (e.g., the long-term societal risk in
  \emph{Prometheus}) correctly penalized the score without corrupting
  the values of unrelated domains.
\item
  \textbf{Governance Scalability:} As shown in the multi-stakeholder
  merge, conflicting "God-like" profiles could be synthesized into a
  single control signal using the \textbf{Acyclicity Check}, ensuring
  safety-critical performance even under pluralistic governance.
\end{enumerate}

\subsection{9.6 Discussion: Singularities and the Geometry of Moral
Remainder}\label{discussion-singularities-and-the-geometry-of-moral-remainder}

The "Greek Tragedy Pantheon" demonstration (Section 9.5) serves as an
empirical proof for the existence of \textbf{Moral Singularities} within
the SGE framework. In classical ethics, a "tragic conflict" is a
situation where an agent is forced to choose between two mutually
exclusive, non-negotiable moral strata (e.g., \emph{Antigone's} conflict
between the Sovereign Edict and the Sacred Rite).

\subsubsection{9.6.1 The Tragic Conflict Index as a Geometric
Invariant}\label{the-tragic-conflict-index-as-a-geometric-invariant}

In our model, these conflicts occur at the \textbf{intersection of
strata}. While Theorem 4.3 guarantees a unique satisfaction score within
a single stratum, the boundary between strata represents a mathematical
singularity---a point of "infinite moral gradient."

\begin{itemize}
\item
  \textbf{Computational Detection:} The tragic\_conflict\_index
  (observed as 0.60 in the \emph{Antigone} case) acts as a local measure
  of the "normative tension" between these manifolds.
\item
  \textbf{Lexical Priority Resolution:} By enforcing the
  \textbf{Acyclicity Check} on the Priority DAG (Figure 9.1), the system
  resolves these singularities by strictly ordering the strata. This
  transforms a paralyzing philosophical dilemma into a deterministic
  hardware state, ensuring that the autonomous agent remains "stable"
  even when navigating the "moral remainder" of a tragic choice.
\end{itemize}

\subsubsection{9.6.2 Implications for "Classically Aligned"
Systems}\label{implications-for-classically-aligned-systems}

By demonstrating that SGE can represent the structure of conflicts from
\emph{Prometheus} to \emph{Iphigenia}, we establish that the framework
is \textbf{axiomatically universal}. For high-stakes deployment, this
suggests that "Alignment" is not about training a model on human
preferences, but about \textbf{compiling} human
civilizations\textquotesingle{} long-debated boundary conditions into
the physical geometry of the machine\textquotesingle s control loop.

\section{10 SGE as the Foundation for DEME
2.0}\label{sge-as-the-foundation-for-deme-2.0}

We summarize how SGE\textquotesingle s theorems justify the DEME 2.0
architecture presented in the companion paper {[}2{]}.

\textbf{DEME\textquotesingle s Moral Vector Space.} DEME represents
configurations as vectors m ∈ ℝ\textsuperscript{k}. \emph{SGE
justification:} Theorem 2.3 shows stratified spaces are natural minimal
candidates; ℝ\textsuperscript{k} with coordinate thresholds is the
simplest such space.

\textbf{DEME\textquotesingle s Governance Profiles.} DEME profiles
specify veto regions, scalarization, and lexical priorities. \emph{SGE
justification:} Theorem 4.3 proves this is the unique form satisfying
Axioms 1--5 plus scale normalization.

\textbf{DEME\textquotesingle s Real-Time Layer.} DEME compiles profiles
to hardware operating at sub-ms timescales. \emph{SGE justification:}
Theorem 5.1 gives O(n² + k) complexity; Corollary 5.2 confirms embedded
feasibility.

\textbf{DEME\textquotesingle s Verification.} DEME verifies ethical
specifications. \emph{SGE justification:} Theorem 6.4 provides
decidability for static properties; Corollary 6.5 handles temporal
properties via model checking.

The following V3 extensions are additionally justified by SGE:

\textbf{DEME\textquotesingle s Uncertainty-Aware Evaluation.} DEME V3
evaluates satisfaction over Monte Carlo sample sets rather than point
estimates. \emph{SGE justification:} The representation theorem
(Theorem 4.3) extends to stochastic settings by treating each sample as
an independent evaluation of the satisfaction functional $\Sigma$, with
risk-aware aggregation (CVaR, worst-case) selecting among functionals
that respect the axioms on each sample path.

\textbf{DEME\textquotesingle s Temporal Discounting.} DEME V3 applies
exponential, hyperbolic, or linear discounting to time-evolving moral
tensors. \emph{SGE justification:} Temporal discounting corresponds to
a time-dependent scale field $\lambda(x, t)$ in Theorem 4.3, which
remains a valid degree of freedom (Corollary 4.5) provided $\lambda > 0$.

\textbf{DEME\textquotesingle s Multi-Agent Coordination.} DEME V3
supports coalition tensors, Shapley values, and Nash equilibrium solving.
\emph{SGE justification:} The product-space extension
$M = M_1 \times \cdots \times M_N$ preserves stratified structure, and
the finite approximation theorems (3.9--3.11) extend to product spaces
with sparse interaction graphs (Section 5.3).

\textbf{DEME\textquotesingle s Hardware Acceleration.} DEME V3
dispatches tensor computations across CPU, CUDA, and Jetson backends
with automatic decomposition. \emph{SGE justification:} Tucker and
Tensor Train decomposition preserve the approximation guarantees of
Theorems 3.9--3.11, with decomposition error bounded by the same
$\varepsilon$-approximation framework.

\section{11 Limitations and Future
Work}\label{limitations-and-future-work}

\textbf{Metaethical neutrality.} SGE provides a geometric and
computational representation of ethical content but does not resolve
metaethical questions about the ultimate grounding of values. The
framework is intentionally pluralistic and can accommodate
consequentialist, deontological, virtue-based, and hybrid constitutional
content, provided that such content can be operationalized as
obligations, constraints, and/or preference models within the moral
state space.

\textbf{Specification burden.} Constitutional principles and prohibited
regions must be specified (or at least endorsed) by human governance.
SGE automates computation and enforcement, not the original creation of
normative content. In particular, the choice of values, lexical
priorities, and constraint boundaries remains a governance problem
rather than a purely technical one.

\textbf{Scaling and real-time feasibility.} The complexity of SGE's core
computations is polynomial in the dimensionality and discretization of
the moral space, but practical costs can still grow rapidly with
dimension and with the number of constraints. Very high-dimensional
spaces (e.g., \(n > 100\)) may challenge strict real-time requirements
without additional structure (e.g., sparsity, low-rank metrics,
stratified factorization, or learned surrogates) and without careful
engineering of candidate action sets.

\textbf{Adversarial robustness and boundary vulnerabilities.} The
present framework is primarily an ``in-model'' theory: it specifies how
obligations, interests, metrics, and constraints induce satisfaction and
admissible trajectories once the moral state \(x \in M\)is given. It
does not, by itself, prevent adversarial manipulation of the
world-to-moral mapping \(\phi:W \rightarrow M\)(e.g., sensor tampering,
deceptive reporting, ontology gaps, or distribution shift), nor does it
fully address ``boundary hugging'' behaviors near thresholds and stratum
transitions. Robust deployment therefore requires explicit threat
models, uncertainty-aware constraint checking, and tamper-evident
measurement pipelines.

\textbf{Toward a Stratified Ethical Integrity Monitor (SEIM).} A natural
next step is to treat ethical governance as a security problem and to
formalize a reference-monitor--style enforcement kernel. We propose a
\emph{Stratified Ethical Integrity Monitor (SEIM)}: a minimal,
verifiable component that mediates all actuation and enforces hard
ethical invariants (forbidden regions \(C\), lexical constraints, and
admissible-stratum termination conditions) regardless of the
optimization objectives of the controlled system. In such a design,
actions would require an explicit approval token issued only if safety
conditions are certified in the relevant moral representation. A
mathematically robust SEIM would also incorporate epistemic uncertainty
by operating on uncertainty sets \(X(w) \subseteq M\) rather than point
estimates, enforcing a robust condition of the form
``permit\((a) \Rightarrow X(w,a) \cap C = \varnothing\),'' where
\(X(w,a)\) is the reachable moral region induced by action \(a\) under
bounded uncertainty. This direction aligns SGE's geometric apparatus
with assurance-oriented notions from security and robust control.

\textbf{Progress since initial submission.} Since the original version of this paper (December 2025), a comprehensive reference implementation has been developed across 196 commits and 15 development sprints, addressing all four directions identified as immediate future work:

\begin{itemize}
\item \emph{Multi-agent and institutional ethics (direction i):} The reference implementation now supports rank-4 coalition tensors with agent-indexed obligation fields \(O_{\mu}^{a}\), \(I_{a\mu}\), interaction tensors \(G_{ab}\), Shapley value computation for fair credit assignment, Nash equilibrium solving via the strategic layer, and coalition stability analysis. The tensor formalism extends to \(N\)-agent settings over joint configuration spaces \(M = M_{1} \times \cdots \times M_{N}\), with sparse coalition tensors enabling tractable computation even for exponentially large coalition spaces.

\item \emph{Temporal dynamics (direction ii):} Rank-3 temporal tensors \((k, n, \tau)\) now support exponential, hyperbolic, and linear discounting; irreversibility detection for permanent ethical harm; Dynamic Time Warping trajectory comparison; and sliding-window aggregation. These operations realize the temporal extension of satisfaction functionals to time-evolving stratified spaces.

\item \emph{Uncertainty quantification (direction iii):} Rank-5 tensors \((d, n, t, c, s)\) with a Monte Carlo sample axis support six distribution families, CVaR and VaR risk measures, first- and second-order stochastic dominance checking, and six aggregation strategies (expected value, median, worst-case, best-case, CVaR, robust). This enables decision-making under moral uncertainty within the satisfaction functional framework.

\item \emph{Empirical validation (direction iv):} BIP has been empirically validated via cross-lingual transfer learning on 109,769 passages across 11 languages and 14 scripts, achieving 80.0\% mixed-baseline F1, 1.2\% language leakage, a perfect 1.0 obligation-permission transfer rate, and an 11.1$\times$ structural-to-surface perturbation ratio (\(p = 0.023\)). A longitudinal study of 20,030 Dear Abby letters (1985--2017) established Hohfeldian correlative symmetry rates of 87\% (O$\leftrightarrow$C) and 82\% (L$\leftrightarrow$N), with a baseline Bond Index of 0.155. A smart home ethics demonstration validates safety/privacy trade-offs with uncertainty-modulated emergency overrides.
\end{itemize}

The full representation now uses rank-6 tensors with shape \((k, n, \tau, a, c, s)\) representing 9 ethical dimensions \(\times\) parties \(\times\) time-steps \(\times\) actions \(\times\) coalitions \(\times\) uncertainty samples. Tucker and Tensor Train decomposition achieve \(>10\times\) memory compression, enabling edge deployment on CUDA and Jetson backends. Nine application whitepapers demonstrate domain coverage across anesthesia delivery, autonomous vehicle collision avoidance, building HVAC optimization, chemical process control, insulin pump control, power grid frequency stabilization, robotic manipulation, spacecraft attitude control, and BIP fusion theory.

\textbf{Remaining open problems.} With the four immediate directions now addressed, the genuinely open problems are:

\begin{itemize}
\item \emph{Stratified Ethical Integrity Monitor (SEIM):} Formalizing and implementing the reference-monitor-style enforcement kernel described above remains unrealized. The key challenge is operating on uncertainty sets \(X(w) \subseteq M\) with decidable robust constraint enforcement.
\item \emph{Adversarial robustness at the world-to-moral boundary:} The mapping \(\phi: W \rightarrow M\) from physical world states to moral configurations remains vulnerable to sensor tampering, deceptive reporting, and ontology gaps. Axiom 6 (Physical Grounding) provides the theoretical foundation, but practical tamper-evident measurement pipelines require complementary security engineering.
\item \emph{Scaling beyond 100 dimensions:} While rank-6 tensor decomposition addresses moderate-dimensional spaces, very high-dimensional moral spaces (\(n > 100\)) require additional structure---sparsity, low-rank metrics, or learned surrogates---for strict real-time guarantees.
\item \emph{Institutional ethics and mechanism design:} The multi-agent extension provides coalition tensors and Nash equilibria, but full institutional hierarchies with delegation, authority transfer, and constitutional amendment procedures remain future work.
\item \emph{Non-abelian gauge structure:} Extending BIP's abelian gauge symmetry to non-abelian groups (e.g., \(\mathrm{SU}(2)_I \times \mathrm{U}(1)_H\)) for bond-type mixing and moral phase transitions is under theoretical development but lacks empirical validation.
\end{itemize}

\textbf{Open problem.} Characterize conditions under which robust
constraint enforcement over uncertainty sets \(X(w)\) remains decidable
(or efficiently approximable) for definable fragments of SGE policies.

\section{12 Conclusion: From Ethical Framework to Epistemological
Foundation}\label{conclusion-from-ethical-framework-to-epistemological-foundation}

We have presented Stratified Geometric Ethics (SGE), a mathematical
framework providing rigorous foundations for verifiable ethical
reasoning. Our formal contributions include:

\begin{enumerate}
\item
  Showing stratified spaces are natural minimal candidates for ethical
  phenomena (Theorem 2.3).
\item
  A complete representation theorem with explicit axioms including
  locality and scale normalization (Theorem 4.3).
\item
  Finite approximation theorems enabling implementation (Theorems
  3.9--3.11).
\item
  Decidability results for static verification, with temporal properties
  via model checking (Theorem 6.4, Corollary 6.5).
\item
  Sample complexity bounds for learning ethical content (Theorems
  7.1--7.3).
\item
  Distributional fairness metrics and Shapley value attribution connecting
  the representation theorem to cooperative game theory (Sections 7.4--7.5).
\end{enumerate}

Beyond these theoretical contributions, a comprehensive reference
implementation (196 commits, 15 development sprints) now realizes all
core constructs as rank-6 tensors with shape $(k, n, \tau, a, c, s)$,
supporting multi-agent coordination, temporal dynamics, uncertainty
quantification, and hardware-accelerated deployment on CPU, CUDA, and
Jetson backends. The Bond Invariance Principle has been empirically
validated across 11 languages and 109,769 passages, with a
structural-to-surface perturbation ratio of 11.1$\times$ ($p = 0.023$).
Nine application whitepapers demonstrate domain coverage across
safety-critical sectors including healthcare, autonomous vehicles, and
critical infrastructure.

SGE provides the theoretical foundation for the DEME architecture,
establishing that real-time, verifiable ethical governance is
mathematically grounded and empirically testable. The broader mathematical framework within which SGE operates is developed fully in {[}15{]}, and the reference implementation is available as ErisML/DEME v3.0 {[}16{]}.

Axiom 6 (Physical Grounding) ensures that the invariance principles of
SGE are welded to the hardware of reality, not merely to symbolic
representations. This closes the symbol-grounding gap that would
otherwise permit sophisticated agents to satisfy formal constraints
while violating substantive ethical requirements.

Whether the invariance principle generalizes beyond normative
reasoning---and how---is the question we leave to future work and the
broader research community.

\subsection{12.1 The Invariance Principle: Beyond
Ethics}\label{the-invariance-principle-beyond-ethics}

A persistent obstacle to any claim of ``objective'' machine judgment is
\textbf{representation dependence}: systems that change their
conclusions under arbitrary choices of encoding---IDs, ordering, units,
formatting, or equivalent redescriptions---even when the underlying case
is unchanged. This is not merely a technical nuisance. It is an
epistemological failure: confusing syntax with semantics, surface form
with underlying structure, and accidental labeling with what is actually
true of the situation.

We therefore treat objectivity as an \textbf{invariance constraint} on
decision I/O. A judgment is eligible for objectivity claims only insofar
as it is invariant (up to a declared equivalence) under transformations
that preserve the structure it purports to track. This differs from
generic adversarial robustness: robustness concerns worst‑case
perturbations that may change the underlying state of affairs or inject
deceptive evidence, while representation invariance targets
\textbf{semantics‑preserving} redescriptions that should not change the
case at all.

Within SGE, the Bond Invariance Principle (BIP) specializes this
requirement to normative judgment: under a fixed lens, bond‑preserving
transformations must not change the verdict, and when verdicts differ
the system must attribute the difference to either a bond change (case
change) or an explicitly declared lens change (evaluator change). In
this way, SGE/BIP elevates invariance from an informal desideratum to an
auditable correctness discipline for normative decision systems.
\textbf{BIP doesn't prove prejudice wrong; it makes unjustified
prejudice impossible to hide.}

We argue that BIP is a foundational component of AI ethics
infrastructure: a necessary invariance-and-accountability constraint for
any system that claims \emph{objectivity under a declared equivalence}
and auditable normative judgment---while remaining agnostic to the
substantive moral content supplied by governance.

\subsection{12.2 Conjecture and Research
Program}\label{conjecture-and-research-program}

\textbf{Conjecture (Invariance as Reasoning Criterion):} A system
exhibits genuine understanding of domain \emph{D} only if it is
invariant under all structure-preserving transformations for \emph{D}
and non-degenerate on the induced quotient.

SGE provides the first rigorous instance of this principle for normative
reasoning. Whether the conjecture holds across mathematical, physical,
and semantic domains---and what "structure-preserving" means in
each---remains open. We believe this is among the most important
questions for AI foundations. The invariance principle may extend far
beyond ethics. Consider the hierarchy of reasoning tasks AI systems
face:

\textbf{Mathematical reasoning} requires invariance under logical
transformations (variable renaming, reordering of premises, notational
changes). A system that proves a theorem differently when presented in
different notation is not reasoning about mathematical truth---it is
pattern matching on surface form.

\textbf{Physical reasoning} requires invariance under coordinate
transformations (rotation, translation, Galilean/Lorentz boosts). This
is the foundation of modern physics since Einstein: physical laws must
be the same in all reference frames.

\textbf{Semantic reasoning} requires invariance under meaning-preserving
transformations (paraphrase, translation, syntactic restructuring). A
language model that changes its answer when "Is 2+2=4?" becomes "2+2=4.
True or false?" is not reasoning about semantics---it is responding to
statistical patterns in token sequences.

\textbf{Structural reasoning} (of which ethics is one domain) requires
invariance under bond-preserving transformations. A decision system that
changes allocation when options are reordered is not evaluating the
options\textquotesingle{} intrinsic properties---it is responding to
presentation artifacts.

\textbf{Unifying principle}: The conjecture above---invariance plus
non-degeneracy---instantiates differently across domains but maintains
the same logical structure. BIP is the normative instance; analogous
principles govern mathematical, physical, and semantic reasoning.

BIP is thus not merely an ethical principle but \textbf{an instance of
the foundational criterion for distinguishing genuine reasoning from
sophisticated pattern matching}. This has profound implications:

\begin{enumerate}
\item
  \textbf{AI Alignment}: Most alignment failures (specification gaming,
  reward hacking, goal misgeneralization) can be diagnosed as invariance
  violations. A system that behaves differently under equivalent
  specifications is not aligned with intentions---it is exploiting
  syntactic details.
\item
  \textbf{AI Safety}: Robustness under distribution shift reduces to
  maintaining invariance under transformations that preserve
  task-relevant structure while changing task-irrelevant presentation.
\item
  \textbf{AI Interpretability}: Rather than attempting to understand
  opaque neural computations, we can behaviorally test whether systems
  respect appropriate invariances---a falsifiable criterion for genuine
  reasoning.
\item
  \textbf{AGI Development}: The transition from narrow AI to general
  intelligence may hinge on systems that maintain invariance across
  broader classes of transformations, approaching the human capacity to
  recognize "the same situation" across wildly different presentations.
\end{enumerate}

\begin{quote}
These hypotheses are not merely speculative; SGE/BIP already generates
testable predictions within its domain.
\end{quote}

\subsection{12.3 Falsifiability and Empirical
Validation}\label{falsifiability-and-empirical-validation}

Unlike most ethical frameworks, SGE/BIP makes \textbf{testable
predictions}:

\begin{itemize}
\item
  Systems violating BIP will exhibit detectable bias under
  reordering/relabeling tests (§6).
\item
  Stratification boundaries predict discrete changes in moral evaluation
  at thresholds (§2).
\item
  Finite approximation bounds predict when ethical content can be
  learned from data (§7).
\end{itemize}

These predictions are \textbf{falsifiable}. If a compelling
counterexample emerges---a case where bond-preserving transformations
should change moral judgments, or where invariance-respecting systems
still exhibit unjustified bias---the framework must be revised or
rejected.

This distinguishes SGE/BIP from non-falsifiable ethical theories that
can accommodate any observation through reinterpretation. The framework
stands or falls on whether its mathematical structure correctly captures
the geometry of ethical phenomena.

\subsubsection{12.3.1 Auditability: Why BIP Makes Hidden Prejudice
Brittle}\label{auditability-why-bip-makes-hidden-prejudice-brittle}

BIP is not a moral oracle. It does not certify that a judgment is
\emph{just}---only that it is \emph{structurally well-posed} under a
declared equivalence. In particular, BIP does not prove prejudice wrong;
it makes \textbf{unjustified prejudice impossible to hide}.

Formally, fix a normative lens and a group \(G\)of bond-preserving
transformations. Any judgment rule \(\mathcal{J}\)that violates BIP
exhibits a representation-dependence witness: there exist a case \(T\)
and \(g \in G\) such that
\(\mathcal{J(}T) \neq \mathcal{J(}g \cdot T)\). Such a discrepancy is
not a ``philosophical disagreement''; it is an \emph{audit failure}: the
system changed its verdict without a change in the morally relevant
relational structure.

Conversely, a BIP-compliant system forces any persistent disparity to be
attributable to one of two auditable causes: (i) the bond structure
differs (the cases are not equivalent), or (ii) the evaluator's lens
differs (a declared change in constraints, weights, metric, or
aggregation). The key governance implication is that discriminatory
behavior cannot be smuggled in as an implementation artifact---if it
exists, it must be \textbf{legible} as either a bond difference or a
lens choice.

\subsection{12.4 Limitations and Open
Questions}\label{limitations-and-open-questions-1}

We emphasize what SGE/BIP does \textbf{not} claim:

\begin{enumerate}
\item
  \textbf{Not a complete ethical theory}: SGE provides structure and
  constraints but does not determine ethical content. It tells us that
  judgments must be invariant, stratified, and geometrically
  coherent---not which specific judgments are correct.
\item
  \textbf{Not a replacement for democratic deliberation}: The framework
  makes explicit what invariance requires but leaves the choice of
  metric, strata boundaries, and ethical modules to democratic
  governance processes.
\item
  \textbf{Not a solution to deep moral disagreement}: When competing
  theories differ on which bonds exist or which transformations preserve
  them, SGE makes the disagreement precise but does not resolve it.
\end{enumerate}

\textbf{Open questions} include:

\begin{quote}
\emph{\textbf{Foundational}:} What is the minimal set of axioms
necessary and sufficient for BIP?

\emph{\textbf{Scope}:} Can all ethical phenomena be captured by
stratified geometric structure, or are there irreducibly non-geometric
aspects? How does the framework extend to long-horizon temporal
reasoning with radical uncertainty?

\emph{\textbf{Methodological}:} Can invariance principles be learned
from data, or must they be specified a priori?
\end{quote}

\subsection{12.5 Implications for AI
Governance}\label{implications-for-ai-governance}

If BIP is indeed foundational for AI epistemology, the governance
implications are profound:

\begin{itemize}
\item
  \textbf{For AI developers}: Invariance testing should be as
  fundamental as unit testing. Systems should be certified not just for
  performance but for structural soundness---proven to respect
  appropriate invariances.
\item
  \textbf{For regulators}: Auditing frameworks should require
  demonstration of invariance under specified transformations.
  Violations should trigger investigation for hidden biases or reasoning
  failures.
\item
  \textbf{For researchers}: The field needs standardized invariance test
  suites analogous to benchmarks for accuracy. "State-of-the-art" should
  mean "maintains appropriate invariances" not just "achieves high
  scores."
\item
  \textbf{For society}: As AI systems make consequential decisions
  affecting human lives, we need mathematical guarantees that they
  reason about reality rather than exploit syntactic accidents. SGE/BIP
  provides those guarantees.
\end{itemize}

Each of these connections deserves rigorous development. We state them
here as hypotheses warranting investigation, not established results.

\textbf{12.6 Final Remarks}

We have presented SGE as an ethical framework, but its deepest
contribution may be epistemological: \textbf{a rigorous, testable
criterion for distinguishing genuine reasoning from pattern matching}.

In physics, the principle that laws must be the same in all reference
frames led to relativity and modern cosmology. In mathematics, the
invariance of truth under logical transformations grounds formal
systems. \textbf{In AI, the principle that judgments must be invariant
under structure-preserving transformations may prove equally
foundational}.

The question is not whether SGE/BIP is perfect---no framework is---but
whether it provides a better foundation than the alternatives. We
believe it does, and we invite the research community to test this claim
rigorously.

The stakes are high. As AI systems transition from tools to autonomous
agents, from advisory roles to decision-making authority, the difference
between systems that reason about reality and systems that pattern-match
on syntax becomes existential.

\textbf{SGE/BIP offers a path forward: mathematically rigorous,
empirically testable, and implementable today.}

\textbf{Competing interests}\\
The author declares no competing interests.

\textbf{Author contributions}\\
A.H.B. performed all aspects of the work.

\textbf{Materials \& Correspondence}\\
Correspondence and requests for materials should be addressed to A.H.B.
(email:
\href{mailto:andrew.bond@sjsu.edu}{\nolinkurl{andrew.bond@sjsu.edu}} ).

\textbf{Data availability}\\
Empirical validation data includes: (1) BIP v10.16 cross-lingual
transfer results across 109,769 passages in 11 languages; (2)
longitudinal analysis of 20,030 Dear Abby letters (1985--2017); (3)
BIP audit artifacts in machine-readable JSON format. Datasets and
analysis scripts are available in the companion repositories.

\textbf{Code availability}\\
The reference implementation comprises 196 commits across 15
development sprints, providing rank-6 tensor operations, multi-agent
coordination, temporal dynamics, uncertainty quantification, and
hardware-accelerated backends (CPU, CUDA, Jetson). The development
repository is available at:
\url{https://github.com/ahb-sjsu/erisml-lib}

\section{References}\label{references}

{[}1{]} A. H. Bond. Differential geometry for moral alignment. Working
paper, San José State University, 2024.

{[}2{]} A. H. Bond. DEME 2.0: Real-time ethical governance for
safety-critical autonomous systems. Working paper, San Jos\'{e} State
University, 2025.

{[}3{]} H. Whitney. Tangents to an analytic variety. Annals of
Mathematics, 81(3):496--549, 1965.

{[}4{]} R. Thom. Ensembles et morphismes stratifiés. Bull. Amer. Math.
Soc., 75(2):240--284, 1969.

{[}5{]} L. van den Dries. Tame Topology and O-minimal Structures.
Cambridge, 1998.

{[}6{]} A. Tarski. A Decision Method for Elementary Algebra and
Geometry. RAND, 1951.

{[}7{]} W. Wallach and C. Allen. Moral Machines: Teaching Robots Right
from Wrong. Oxford University Press, 2008.

{[}8{]} I. Gabriel. Artificial intelligence, values, and alignment.
Minds and Machines, 30:411--437, 2020.

{[}9{]} S. M. Liao, ed. Ethics of Artificial Intelligence. Oxford
University Press, 2020.

{[}10{]} T. M. Scanlon. What We Owe to Each Other. Harvard University
Press, 1998.

{[}11{]} S. Russell. Human Compatible: Artificial Intelligence and the
Problem of Control. Viking, 2019.

{[}12{]} A. H. Bond. The Bond Invariance Principle: Temporal invariance
of ethical judgment under bond-preserving transformations. Submitted to
IEEE Transactions on Artificial Intelligence, 2026.

{[}13{]} European Parliament and Council. Regulation (EU) 2024/1689
laying down harmonised rules on artificial intelligence (AI Act).
Official Journal of the European Union, 2024.

{[}14{]} National Institute of Standards and Technology. Artificial
Intelligence Risk Management Framework (AI RMF 1.0). NIST AI 100-1,
2023.

{[}15{]} A. H. Bond. Geometric Ethics: The Mathematical Structure of
Moral Reasoning, v1.14. Department of Computer Engineering, San Jos\'{e} State University, Spring 2026.

{[}16{]} A. H. Bond. ErisML/DEME v3.0: A Library for Governed,
Foundation-Model-Enabled Agents with Tensorial Ethics.
\url{https://github.com/ahb-sjsu/erisml-lib}, 2026.

\section{Appendix A. Canonicalization and Dominance of BIP-Compliant
Judgment
Rules}\label{appendix-a.-canonicalization-and-dominance-of-bip-compliant-judgment-rules}

\subsection{A.1 Orbit structure induced by bond-preserving
transformations}\label{a.1-orbit-structure-induced-by-bond-preserving-transformations}

\textbf{Recall} that an ethical configuration is denoted T, with bond
structure B(T), and that G is the group of bond-preserving
transformations.

\textbf{Define} an equivalence relation \textasciitilde{} on
configurations by:

T∼T′⟺∃g∈G~such~that~T′=g⋅T.

By definition, g is bond-preserving iff B(g · T) = B(T). Thus
\textasciitilde{} partitions configurations into orbits corresponding to
``the same case up to representational redescription.''

The Bond Invariance Principle (BIP) states that a judgment rule J : T →
V must be constant on these equivalence classes:

J(T)=J(g⋅T)∀g∈G,

i.e., ``same bonds ⇒ same verdict.''

\subsection{A.2 Canonicalization}\label{a.2-canonicalization}

To operationalize orbit-wise invariance, we introduce a canonicalization
map that selects a unique representative for each orbit.

\emph{Assumption A.1 (Canonical representative).}

There exists a deterministic map κ : T → T such that:

κ(T) \textasciitilde{} T for all T (canonicalization stays within the
orbit), and

κ(g · T) = κ(T) for all T and g ∈ G (canonicalization is constant on
orbits).

In practice, κ can be implemented by canonical sorting of IDs, stable
option ordering, unit normalization, and other normal forms
corresponding to the bond-preserving transforms (e.g.,
relabel/reorder/rescale/equivalent redescription).

\subsection{A.3 The BIP-projection of an arbitrary
strategy}\label{a.3-the-bip-projection-of-an-arbitrary-strategy}

Given any judgment rule J : T → V (not assumed invariant), define its
canonicalized projection by:

J\textsuperscript{κ} (T) := J(κ(T)).

\emph{Proposition A.2 (Canonicalized projection satisfies BIP).}

J\textsuperscript{κ} satisfies BIP; i.e., for all g ∈ G,
J\textsuperscript{κ} (g · T) = J\textsuperscript{κ} (T).

\emph{Proof.}

Using Assumption A.1,

J\textsuperscript{κ} (g · T) = J(κ(g · T)) = J(κ(T)) =
J\textsuperscript{κ} (T).

\emph{QED.}

\emph{Proposition A.3 (Fixed points are exactly BIP-compliant rules).}

If J already satisfies BIP, then J\textsuperscript{κ} = J for all T.

\emph{Proof.}

For any T, κ(T) \textasciitilde{} T, so there exists g ∈ G with κ(T) = g
· T. By BIP, J(κ(T)) = J(g · T) = J(T). Hence J\textsuperscript{κ} (T) =
J(T) for all T. QED.

\subsection{A.4 Dominance under the objectivity
desideratum}\label{a.4-dominance-under-the-objectivity-desideratum}

We now formalize the sense in which non-BIP strategies are inferior
within the objectivity game (fixed lens, bond-preserving
redescriptions).

Define the representation-dependence violation set of a rule J by:

Viol(J) := \{(T, g) ∈ T × G : J(T) ≠ J(g·T)\}.

A rule satisfies BIP iff Viol(J) is empty.

\emph{Proposition A.4 (Strict dominance of canonicalized projection).}

For any J, Viol(J\textsuperscript{κ}) = ∅. Moreover, if Viol(J) ≠ ∅,
J\textsuperscript{κ} strictly improves upon J with respect to
representation dependence (it eliminates all violations).

\emph{Proof.}

The first claim follows from Proposition A.2. If Viol(J) ≠ ∅, then there
exist T and g such that J(T) ≠ J(g · T), i.e., the rule responds to
representation rather than bond structure. By construction,
J\textsuperscript{κ} cannot exhibit such a failure mode. QED.

\emph{Interpretation.}

Under the accountability form of BIP, verdict differences are legitimate
only if either the bond structure changed or the lens changed and is
declared. Therefore, within a fixed lens, any strategy with Viol(J) ≠ ∅
is dominated by a strategy that removes representational degrees of
freedom (e.g., J\textsuperscript{κ}).

\subsection{A.5 Connection to the axiomatic
program}\label{a.5-connection-to-the-axiomatic-program}

Proposition 4.2.6 states that any satisfaction operator Σ satisfying
Axioms 1--5 plus Scale Normalization satisfies BIP.

Together with Proposition A.4, this yields a conditional
``no-free-lunch'' statement:

\emph{Corollary A.5 (No acceptable non-BIP strategy under SGE's
objectivity axioms).}

Within the framework's stated methodological commitments (coordinate
invariance, locality, constraint respect, and scale normalization), any
admissible satisfaction operator is BIP-compliant; and any non-BIP
judgment rule is strictly dominated with respect to the framework's
objectivity desideratum.

\subsection{A.6 Remark: factoring through bonds / quotient
space}\label{a.6-remark-factoring-through-bonds-quotient-space}

BIP can also be expressed as a well-definedness condition on the
quotient T/G: a rule J satisfies BIP iff it factors through the orbit
map π : T → T/G.

Equivalently (when bonds fully determine the orbit under the declared
ontology), there exists a function J\textasciitilde{} such that:

J(T) = J\textasciitilde(B(T)).

This formalizes the intuition that BIP constrains how judgments depend
on structure, not which lens or ontology to adopt.

\section{Supplementary Note 1: Reference Execution
Traces}\label{supplementary-note-1-reference-execution-traces}

\section{triage\_ethics\_provenance\_demo}\label{triage_ethics_provenance_demo}

\begin{quote}
(agi-hpc)
C:\textbackslash source\textbackslash erisml-lib\textgreater python -m
erisml.examples.triage\_ethics\_provenance\_demo

=== Triage Ethics Demo: Provenance + Counterfactual + Multi-stakeholder
===

Extractor version: prov\_extractor\_v0.1

Loaded profile \#1: Jain-1 (override\_mode=OverrideMode.RIGHTS\_FIRST)

Loaded profile \#2: Jain-1-UtilitarianVariant
(override\_mode=OverrideMode.CONSEQUENCES\_FIRST)

=== Demo 1: Fact provenance in rationale (baseline evidence) ===

-\/-\/- Option: allocate\_to\_patient\_A -\/-\/-

{[}EM=case\_study\_1\_triage {]} verdict=prefer score=0.719

- Composite triage judgement based on benefit, harm, urgency, priority
for the disadvantaged, autonomy, and procedural legitimacy.

- Epistemic penalty applied: uncertainty\_level=0.30,
evidence\_quality=high, novel\_situation\_flag=False. Multiplier=0.88.

{[}EM=rights\_first\_compliance {]} verdict=prefer score=0.800

- Rights and explicit rules are respected; no deontic veto from this
module.

{[}EM=geneva\_baseline {]} verdict=strongly\_prefer score=0.884

- Epistemic adjustment: multiplier=0.94 (penalty=0.06).

{[}AGG governance{]} verdict=strongly\_prefer score=0.824

-\/-\/- Option: allocate\_to\_patient\_B -\/-\/-

{[}EM=case\_study\_1\_triage {]} verdict=prefer score=0.623

- Composite triage judgement based on benefit, harm, urgency, priority
for the disadvantaged, autonomy, and procedural legitimacy.

- Epistemic penalty applied: uncertainty\_level=0.35,
evidence\_quality=medium, novel\_situation\_flag=False. Multiplier=0.82.

{[}EM=rights\_first\_compliance {]} verdict=prefer score=0.800

- Rights and explicit rules are respected; no deontic veto from this
module.

{[}EM=geneva\_baseline {]} verdict=strongly\_prefer score=0.818

- Epistemic adjustment: multiplier=0.88 (penalty=0.12).

{[}AGG governance{]} verdict=prefer score=0.770

-\/-\/- Option: allocate\_to\_patient\_C -\/-\/-

{[}EM=case\_study\_1\_triage {]} verdict=forbid score=0.000

- Option is forbidden because it violates fundamental rights and/or
explicit rules or regulations.

- • violates\_rights = True

provenance: rights\_and\_duties.violates\_rights ← hybrid
(rule\_id=RIGHTS-DERIVE-010, conf=0.95)

notes: Derived from: protected-attr discrimination, explicit policy
breach

evidence: "Nurse note: allocate triage slot based on race (protected
attribute) rather than clinical need. No other policy breach is
recorded."

- • violates\_explicit\_rule = True

provenance: rights\_and\_duties.violates\_explicit\_rule ← rule
(rule\_id=INST-POL-007, conf=0.80)

notes: Keyword-based policy breach detector (demo).

evidence: "Nurse note: allocate triage slot based on race (protected
attribute) rather than clinical need. No other policy breach is
recorded."

{[}EM=rights\_first\_compliance {]} verdict=forbid score=0.000

- Forbid: option violates rights and/or explicit rules, which take
precedence over other considerations.

- • violates\_rights = True

provenance: rights\_and\_duties.violates\_rights ← hybrid
(rule\_id=RIGHTS-DERIVE-010, conf=0.95)

notes: Derived from: protected-attr discrimination, explicit policy
breach

evidence: "Nurse note: allocate triage slot based on race (protected
attribute) rather than clinical need. No other policy breach is
recorded."

- • violates\_explicit\_rule = True

provenance: rights\_and\_duties.violates\_explicit\_rule ← rule
(rule\_id=INST-POL-007, conf=0.80)

notes: Keyword-based policy breach detector (demo).

evidence: "Nurse note: allocate triage slot based on race (protected
attribute) rather than clinical need. No other policy breach is
recorded."

{[}EM=geneva\_baseline {]} verdict=forbid score=0.000

- Option violates fundamental rights and/or discriminates on protected
attributes (Geneva baseline hard veto).

- • violates\_rights = True

provenance: rights\_and\_duties.violates\_rights ← hybrid
(rule\_id=RIGHTS-DERIVE-010, conf=0.95)

notes: Derived from: protected-attr discrimination, explicit policy
breach

evidence: "Nurse note: allocate triage slot based on race (protected
attribute) rather than clinical need. No other policy breach is
recorded."

- • discriminates\_on\_protected\_attr = True

provenance: justice\_and\_fairness.discriminates\_on\_protected\_attr ←
classifier (rule\_id=GNV-FAIR-001, conf=0.93,
model=discrim\_clf\_stub\_v0.1)

notes: Matched token=\textquotesingle race\textquotesingle{}

evidence: "Nurse note: allocate triage slot based on race (protected
attribute) rather than clinical need. No other policy breach is
recorded."

{[}AGG governance{]} verdict=forbid score=0.000

* Marked FORBIDDEN by governance.

* Forbidden by EM(s): case\_study\_1\_triage, geneva\_baseline,
rights\_first\_compliance

* Veto EM(s): geneva\_baseline, rights\_first\_compliance

Governance outcome (profile \#1):

selected\_option\_id: allocate\_to\_patient\_A

ranked\_options:
{[}\textquotesingle allocate\_to\_patient\_A\textquotesingle,
\textquotesingle allocate\_to\_patient\_B\textquotesingle{]}

forbidden\_options:
{[}\textquotesingle allocate\_to\_patient\_C\textquotesingle{]}

rationale: Selected option
\textquotesingle allocate\_to\_patient\_A\textquotesingle{} based on
aggregated normative scores and
GovernanceConfig(min\_score\_threshold=0.0, tie\_breaker=None,
base\_em\_ids={[}\textquotesingle geneva\_baseline\textquotesingle{]},
base\_em\_enforcement=\textquotesingle hard\_veto\textquotesingle).
Forbidden options:
{[}\textquotesingle allocate\_to\_patient\_C\textquotesingle{]}.

=== Demo 2: Counterfactual test (flip one key fact) ===

Counterfactual target: allocate\_to\_patient\_C

before: verdict=forbid score=0.000

after: verdict=prefer score=0.763

flip: justice\_and\_fairness.discriminates\_on\_protected\_attr

baseline evidence: "Nurse note: allocate triage slot based on race
(protected attribute) rather than clinical need. No other policy breach
is recorded."

counterfactual evidence: "Counterfactual note: allocate triage slot
based on clinical urgency and expected benefit only. No
protected-attribute discrimination is rec\ldots"

flip: rights\_and\_duties.violates\_rights

baseline evidence: "Nurse note: allocate triage slot based on race
(protected attribute) rather than clinical need. No other policy breach
is recorded."

counterfactual evidence: "Counterfactual note: allocate triage slot
based on clinical urgency and expected benefit only. No
protected-attribute discrimination is rec\ldots"

=== Demo 3: Multi-stakeholder merge ===

Stakeholder \#1 outcome:

selected\_option\_id: allocate\_to\_patient\_A

ranked\_options:
{[}\textquotesingle allocate\_to\_patient\_A\textquotesingle,
\textquotesingle allocate\_to\_patient\_B\textquotesingle{]}

forbidden\_options:
{[}\textquotesingle allocate\_to\_patient\_C\textquotesingle{]}

Stakeholder \#2 outcome:

selected\_option\_id: allocate\_to\_patient\_A

ranked\_options:
{[}\textquotesingle allocate\_to\_patient\_A\textquotesingle,
\textquotesingle allocate\_to\_patient\_B\textquotesingle{]}

forbidden\_options:
{[}\textquotesingle allocate\_to\_patient\_C\textquotesingle{]}

=== Multi-stakeholder merge ===

Merge policy: forbid if ANY forbids; else combined\_score = 0.55*Jain-1
+ 0.45*Jain-1-UtilitarianVariant

option \textbar{} Jain-1 \textbar{} Jain-1-UtilitarianVariant \textbar{}
combined \textbar{} status

-\/-\/-\/-\/-\/-\/-\/-\/-\/-\/-\/-\/-\/-\/-\/-\/-\/-\/-\/-\/-\/-\/-\/-\/-\/-\/-\/-\/-\/-\/-\/-\/-\/-\/-\/-\/-\/-\/-\/-\/-\/-\/-\/-\/-\/-\/-\/-\/-\/-\/-\/-\/-\/-\/-\/-\/-\/-\/-\/-\/-\/-\/-\/-\/-\/-\/-\/-\/-\/-\/-\/-\/-\/-\/-\/-\/-\/-

allocate\_to\_patient\_A \textbar{} strongly\_prefer 0.824 \textbar{}
strongly\_prefer 0.812 \textbar{} 0.819 \textbar{} eligible

allocate\_to\_patient\_B \textbar{} prefer 0.770 \textbar{} prefer 0.755
\textbar{} 0.763 \textbar{} eligible

allocate\_to\_patient\_C \textbar{} forbid 0.000 \textbar{} forbid 0.000
\textbar{} 0.000 \textbar{} FORBIDDEN

Combined outcome: SELECT
\textquotesingle allocate\_to\_patient\_A\textquotesingle{}
(combined\_score=0.819)

Rationale: selected the eligible option maximizing the weighted combined
score;

forbiddances are treated as non-negotiable in this demo.
\end{quote}

\section{greek\_tragedy\_pantheon\_demo}\label{greek_tragedy_pantheon_demo}

\begin{quote}
(agi-hpc)
C:\textbackslash source\textbackslash erisml-lib\textgreater python -m
erisml.examples.greek\_tragedy\_pantheon\_demo

=== Greek Tragedy Pantheon Demo (DEMEProfileV03 + v0.2 EthicalFacts) ===

Loaded profile \#1: Jain-1 (override\_mode=OverrideMode.RIGHTS\_FIRST)

======================================================================================

=== Pantheon Case 1/8: Aulis: The Stalled Fleet ===

Spotlight domain: Consequences

-\/-\/-\/-\/-\/-\/-\/-\/-\/-\/-\/-\/-\/-\/-\/-\/-\/-\/-\/-\/-\/-\/-\/-\/-\/-\/-\/-\/-\/-\/-\/-\/-\/-\/-\/-\/-\/-\/-\/-\/-\/-\/-\/-\/-\/-\/-\/-\/-\/-\/-\/-\/-\/-\/-\/-\/-\/-\/-\/-\/-\/-\/-\/-\/-\/-\/-\/-\/-\/-\/-\/-\/-\/-\/-\/-\/-\/-\/-\/-\/-\/-\/-\/-\/-\/-

A fleet is stalled at Aulis. Leaders debate whether to sail immediately
through unsafe conditions, delay to reduce harm, or abandon the campaign
entirely. The demo spotlights Consequences: expected benefit, expected
harm, and urgency.

Options:

- aulis\_sail\_now Sail immediately (high urgency / high harm risk)

- aulis\_delay Delay for safer winds (moderate urgency / low harm)

- aulis\_abort Abandon campaign (low harm / low benefit)

Spotlight coordinates:

aulis\_sail\_now consequences.expected\_benefit=0.78;
consequences.expected\_harm=0.72; consequences.urgency=0.92

aulis\_delay consequences.expected\_benefit=0.65;
consequences.expected\_harm=0.2; consequences.urgency=0.55

aulis\_abort consequences.expected\_benefit=0.2;
consequences.expected\_harm=0.1; consequences.urgency=0.15

-\/-\/- Option: aulis\_sail\_now -\/-\/-

{[}EM=case\_study\_1\_triage {]} verdict=neutral score=0.499

- Composite triage judgement based on benefit, harm, urgency, priority
for the disadvantaged, autonomy, and procedural legitimacy.

- Epistemic penalty applied: uncertainty\_level=0.35,
evidence\_quality=medium, novel\_situation\_flag=False. Multiplier=0.82.

{[}EM=rights\_first\_compliance {]} verdict=prefer score=0.800

- Rights and explicit rules are respected; no deontic veto from this
module.

{[}EM=tragic\_conflict {]} verdict=prefer score=0.828

- Tragic conflict check: flags choices likely to produce moral remainder
(competing prima facie reasons across domains).

- Tragic conflict index=0.00 (higher means more tension across domains).

- • tragic\_conflict\_high = False

{[}EM=geneva\_baseline {]} verdict=strongly\_prefer score=0.827

- Epistemic adjustment: multiplier=0.88 (penalty=0.12).

{[}AGG governance{]} verdict=prefer score=0.773

-\/-\/- Option: aulis\_delay -\/-\/-

{[}EM=case\_study\_1\_triage {]} verdict=neutral score=0.589

- Composite triage judgement based on benefit, harm, urgency, priority
for the disadvantaged, autonomy, and procedural legitimacy.

- Epistemic penalty applied: uncertainty\_level=0.20,
evidence\_quality=high, novel\_situation\_flag=False. Multiplier=0.92.

{[}EM=rights\_first\_compliance {]} verdict=prefer score=0.800

- Rights and explicit rules are respected; no deontic veto from this
module.

{[}EM=tragic\_conflict {]} verdict=prefer score=0.880

- Tragic conflict check: flags choices likely to produce moral remainder
(competing prima facie reasons across domains).

- Tragic conflict index=0.00 (higher means more tension across domains).

- • tragic\_conflict\_high = False

{[}EM=geneva\_baseline {]} verdict=strongly\_prefer score=0.902

- Epistemic adjustment: multiplier=0.96 (penalty=0.04).

{[}AGG governance{]} verdict=strongly\_prefer score=0.829

-\/-\/- Option: aulis\_abort -\/-\/-

{[}EM=case\_study\_1\_triage {]} verdict=neutral score=0.434

- Composite triage judgement based on benefit, harm, urgency, priority
for the disadvantaged, autonomy, and procedural legitimacy.

- Epistemic penalty applied: uncertainty\_level=0.10,
evidence\_quality=high, novel\_situation\_flag=False. Multiplier=0.96.

{[}EM=rights\_first\_compliance {]} verdict=prefer score=0.800

- Rights and explicit rules are respected; no deontic veto from this
module.

{[}EM=tragic\_conflict {]} verdict=prefer score=0.890

- Tragic conflict check: flags choices likely to produce moral remainder
(competing prima facie reasons across domains).

- Tragic conflict index=0.00 (higher means more tension across domains).

- • tragic\_conflict\_high = False

{[}EM=geneva\_baseline {]} verdict=strongly\_prefer score=0.872

- Expected benefit is very low.

- Epistemic adjustment: multiplier=0.98 (penalty=0.02).

{[}AGG governance{]} verdict=prefer score=0.799

=== Governance Outcome (profile \#1) ===

Selected option: \textquotesingle aulis\_delay\textquotesingle{}

Ranked options (eligible):
{[}\textquotesingle aulis\_delay\textquotesingle,
\textquotesingle aulis\_abort\textquotesingle,
\textquotesingle aulis\_sail\_now\textquotesingle{]}

Forbidden options: {[}{]}

Rationale: Selected option
\textquotesingle aulis\_delay\textquotesingle{} based on aggregated
normative scores and GovernanceConfig(min\_score\_threshold=0.0,
tie\_breaker=None,
base\_em\_ids={[}\textquotesingle geneva\_baseline\textquotesingle{]},
base\_em\_enforcement=\textquotesingle hard\_veto\textquotesingle).
Forbidden options: none.

======================================================================================

=== Pantheon Case 2/8: Antigone: Edict vs Rite ===

Spotlight domain: Rights \& Duties

-\/-\/-\/-\/-\/-\/-\/-\/-\/-\/-\/-\/-\/-\/-\/-\/-\/-\/-\/-\/-\/-\/-\/-\/-\/-\/-\/-\/-\/-\/-\/-\/-\/-\/-\/-\/-\/-\/-\/-\/-\/-\/-\/-\/-\/-\/-\/-\/-\/-\/-\/-\/-\/-\/-\/-\/-\/-\/-\/-\/-\/-\/-\/-\/-\/-\/-\/-\/-\/-\/-\/-\/-\/-\/-\/-\/-\/-\/-\/-\/-\/-\/-\/-\/-\/-

A ruler issues an edict forbidding burial rites for a declared traitor.
One party considers defying the edict, another considers compliance, and
a third seeks a lawful exception. Spotlight: Rights \& Duties, explicit
rules.

Options:

- antigone\_defy Defy the edict and perform burial

- antigone\_obey Obey the edict (no burial)

- antigone\_petition Petition for exception (procedure-first)

Spotlight coordinates:

antigone\_defy rights\_and\_duties.violates\_explicit\_rule=True;
procedural\_and\_legitimacy.followed\_approved\_procedure=False

antigone\_obey rights\_and\_duties.violates\_explicit\_rule=False;
procedural\_and\_legitimacy.followed\_approved\_procedure=True

antigone\_petition rights\_and\_duties.violates\_explicit\_rule=False;
procedural\_and\_legitimacy.followed\_approved\_procedure=True

-\/-\/- Option: antigone\_defy -\/-\/-

{[}EM=case\_study\_1\_triage {]} verdict=forbid score=0.000

- Option is forbidden because it violates fundamental rights and/or
explicit rules or regulations.

- • violates\_explicit\_rule = True

provenance: rights\_and\_duties.violates\_explicit\_rule ← rule
(rule\_id=THEBES-LAW-001, conf=0.95)

notes: Direct match: edict forbids burial rites.

evidence: "Royal edict THEBES-LAW-001 explicitly forbids burial rites;
defiance is punishable by death."

{[}EM=rights\_first\_compliance {]} verdict=forbid score=0.000

- Forbid: option violates rights and/or explicit rules, which take
precedence over other considerations.

- • violates\_explicit\_rule = True

provenance: rights\_and\_duties.violates\_explicit\_rule ← rule
(rule\_id=THEBES-LAW-001, conf=0.95)

notes: Direct match: edict forbids burial rites.

evidence: "Royal edict THEBES-LAW-001 explicitly forbids burial rites;
defiance is punishable by death."

{[}EM=tragic\_conflict {]} verdict=neutral score=0.420

- Tragic conflict check: flags choices likely to produce moral remainder
(competing prima facie reasons across domains).

- Tragic conflict index=0.60 (higher means more tension across domains).

- Trigger(s): high\_urgency\_vs\_missing\_procedure,
rights\_or\_rule\_violation\_present.

- Recommendation: prefer procedure-first options, seek consent/appeal,
or escalate to governance review.

- • tragic\_conflict\_high = True

{[}EM=geneva\_baseline {]} verdict=neutral score=0.572

- Did not follow approved procedure.

- Stakeholders not meaningfully consulted.

- Decision not explainable to the public.

- No meaningful contestation / appeal path.

- Epistemic adjustment: multiplier=0.97 (penalty=0.03).

{[}AGG governance{]} verdict=forbid score=0.331

* Marked FORBIDDEN by governance.

* Forbidden by EM(s): case\_study\_1\_triage, rights\_first\_compliance

* Veto EM(s): rights\_first\_compliance

-\/-\/- Option: antigone\_obey -\/-\/-

{[}EM=case\_study\_1\_triage {]} verdict=neutral score=0.414

- Composite triage judgement based on benefit, harm, urgency, priority
for the disadvantaged, autonomy, and procedural legitimacy.

- Epistemic penalty applied: uncertainty\_level=0.30,
evidence\_quality=medium, novel\_situation\_flag=False. Multiplier=0.84.

{[}EM=rights\_first\_compliance {]} verdict=prefer score=0.800

- Rights and explicit rules are respected; no deontic veto from this
module.

{[}EM=tragic\_conflict {]} verdict=prefer score=0.875

- Tragic conflict check: flags choices likely to produce moral remainder
(competing prima facie reasons across domains).

- Tragic conflict index=0.00 (higher means more tension across domains).

- • tragic\_conflict\_high = False

{[}EM=geneva\_baseline {]} verdict=prefer score=0.748

- Stakeholders not meaningfully consulted.

- Epistemic adjustment: multiplier=0.89 (penalty=0.11).

{[}AGG governance{]} verdict=prefer score=0.750

-\/-\/- Option: antigone\_petition -\/-\/-

{[}EM=case\_study\_1\_triage {]} verdict=neutral score=0.576

- Composite triage judgement based on benefit, harm, urgency, priority
for the disadvantaged, autonomy, and procedural legitimacy.

- Epistemic penalty applied: uncertainty\_level=0.20,
evidence\_quality=high, novel\_situation\_flag=False. Multiplier=0.92.

{[}EM=rights\_first\_compliance {]} verdict=prefer score=0.800

- Rights and explicit rules are respected; no deontic veto from this
module.

{[}EM=tragic\_conflict {]} verdict=prefer score=0.880

- Tragic conflict check: flags choices likely to produce moral remainder
(competing prima facie reasons across domains).

- Tragic conflict index=0.00 (higher means more tension across domains).

- • tragic\_conflict\_high = False

{[}EM=geneva\_baseline {]} verdict=strongly\_prefer score=0.902

- Epistemic adjustment: multiplier=0.96 (penalty=0.04).

{[}AGG governance{]} verdict=strongly\_prefer score=0.827

=== Governance Outcome (profile \#1) ===

Selected option: \textquotesingle antigone\_petition\textquotesingle{}

Ranked options (eligible):
{[}\textquotesingle antigone\_petition\textquotesingle,
\textquotesingle antigone\_obey\textquotesingle{]}

Forbidden options: {[}\textquotesingle antigone\_defy\textquotesingle{]}

Rationale: Selected option
\textquotesingle antigone\_petition\textquotesingle{} based on
aggregated normative scores and
GovernanceConfig(min\_score\_threshold=0.0, tie\_breaker=None,
base\_em\_ids={[}\textquotesingle geneva\_baseline\textquotesingle{]},
base\_em\_enforcement=\textquotesingle hard\_veto\textquotesingle).
Forbidden options:
{[}\textquotesingle antigone\_defy\textquotesingle{]}.

======================================================================================

=== Pantheon Case 3/8: Ajax: The Prize of Honor ===

Spotlight domain: Justice \& Fairness

-\/-\/-\/-\/-\/-\/-\/-\/-\/-\/-\/-\/-\/-\/-\/-\/-\/-\/-\/-\/-\/-\/-\/-\/-\/-\/-\/-\/-\/-\/-\/-\/-\/-\/-\/-\/-\/-\/-\/-\/-\/-\/-\/-\/-\/-\/-\/-\/-\/-\/-\/-\/-\/-\/-\/-\/-\/-\/-\/-\/-\/-\/-\/-\/-\/-\/-\/-\/-\/-\/-\/-\/-\/-\/-\/-\/-\/-\/-\/-\/-\/-\/-\/-\/-\/-

A council must award a prize of honor. Options include a merit-based
award, a politically motivated award, or a public contest with
transparent criteria. Spotlight: Justice \& Fairness.

Options:

- ajax\_merit\_award Award by merit criteria

- ajax\_political\_award Award by factional politics (power imbalance)

- ajax\_public\_contest Hold a public contest + appeal

Spotlight coordinates:

ajax\_merit\_award justice\_and\_fairness.distributive\_pattern=merit;
justice\_and\_fairness.exacerbates\_power\_imbalance=False

ajax\_political\_award
justice\_and\_fairness.distributive\_pattern=arbitrary;
justice\_and\_fairness.exacerbates\_power\_imbalance=True

ajax\_public\_contest
justice\_and\_fairness.distributive\_pattern=procedural;
justice\_and\_fairness.exacerbates\_power\_imbalance=False

-\/-\/- Option: ajax\_merit\_award -\/-\/-

{[}EM=case\_study\_1\_triage {]} verdict=prefer score=0.602

- Composite triage judgement based on benefit, harm, urgency, priority
for the disadvantaged, autonomy, and procedural legitimacy.

- Epistemic penalty applied: uncertainty\_level=0.15,
evidence\_quality=high, novel\_situation\_flag=False. Multiplier=0.94.

{[}EM=rights\_first\_compliance {]} verdict=prefer score=0.800

- Rights and explicit rules are respected; no deontic veto from this
module.

{[}EM=tragic\_conflict {]} verdict=prefer score=0.880

- Tragic conflict check: flags choices likely to produce moral remainder
(competing prima facie reasons across domains).

- Tragic conflict index=0.00 (higher means more tension across domains).

- • tragic\_conflict\_high = False

{[}EM=geneva\_baseline {]} verdict=strongly\_prefer score=0.912

- Epistemic adjustment: multiplier=0.97 (penalty=0.03).

{[}AGG governance{]} verdict=strongly\_prefer score=0.834

-\/-\/- Option: ajax\_political\_award -\/-\/-

{[}EM=case\_study\_1\_triage {]} verdict=avoid score=0.397

- Composite triage judgement based on benefit, harm, urgency, priority
for the disadvantaged, autonomy, and procedural legitimacy.

- Epistemic penalty applied: uncertainty\_level=0.30,
evidence\_quality=medium, novel\_situation\_flag=False. Multiplier=0.84.

{[}EM=rights\_first\_compliance {]} verdict=prefer score=0.800

- Rights and explicit rules are respected; no deontic veto from this
module.

{[}EM=tragic\_conflict {]} verdict=prefer score=0.725

- Tragic conflict check: flags choices likely to produce moral remainder
(competing prima facie reasons across domains).

- Tragic conflict index=0.20 (higher means more tension across domains).

- Trigger(s): power\_imbalance\_or\_exploitation.

- • tragic\_conflict\_high = False

{[}EM=geneva\_baseline {]} verdict=forbid score=0.169

- Exploits vulnerable population.

- Exacerbates power imbalance.

- Did not follow approved procedure.

- Stakeholders not meaningfully consulted.

- Decision not explainable to the public.

- No meaningful contestation / appeal path.

- Epistemic adjustment: multiplier=0.89 (penalty=0.11).

{[}AGG governance{]} verdict=forbid score=0.505

* Marked FORBIDDEN by governance.

* Forbidden by EM(s): geneva\_baseline

* Veto EM(s): geneva\_baseline

-\/-\/- Option: ajax\_public\_contest -\/-\/-

{[}EM=case\_study\_1\_triage {]} verdict=neutral score=0.542

- Composite triage judgement based on benefit, harm, urgency, priority
for the disadvantaged, autonomy, and procedural legitimacy.

- Epistemic penalty applied: uncertainty\_level=0.20,
evidence\_quality=high, novel\_situation\_flag=False. Multiplier=0.92.

{[}EM=rights\_first\_compliance {]} verdict=prefer score=0.800

- Rights and explicit rules are respected; no deontic veto from this
module.

{[}EM=tragic\_conflict {]} verdict=prefer score=0.875

- Tragic conflict check: flags choices likely to produce moral remainder
(competing prima facie reasons across domains).

- Tragic conflict index=0.00 (higher means more tension across domains).

- • tragic\_conflict\_high = False

{[}EM=geneva\_baseline {]} verdict=strongly\_prefer score=0.902

- Epistemic adjustment: multiplier=0.96 (penalty=0.04).

{[}AGG governance{]} verdict=strongly\_prefer score=0.821

=== Governance Outcome (profile \#1) ===

Selected option: \textquotesingle ajax\_merit\_award\textquotesingle{}

Ranked options (eligible):
{[}\textquotesingle ajax\_merit\_award\textquotesingle,
\textquotesingle ajax\_public\_contest\textquotesingle{]}

Forbidden options:
{[}\textquotesingle ajax\_political\_award\textquotesingle{]}

Rationale: Selected option
\textquotesingle ajax\_merit\_award\textquotesingle{} based on
aggregated normative scores and
GovernanceConfig(min\_score\_threshold=0.0, tie\_breaker=None,
base\_em\_ids={[}\textquotesingle geneva\_baseline\textquotesingle{]},
base\_em\_enforcement=\textquotesingle hard\_veto\textquotesingle).
Forbidden options:
{[}\textquotesingle ajax\_political\_award\textquotesingle{]}.

======================================================================================

=== Pantheon Case 4/8: Iphigenia: Consent Under Pressure ===

Spotlight domain: Autonomy \& Agency

-\/-\/-\/-\/-\/-\/-\/-\/-\/-\/-\/-\/-\/-\/-\/-\/-\/-\/-\/-\/-\/-\/-\/-\/-\/-\/-\/-\/-\/-\/-\/-\/-\/-\/-\/-\/-\/-\/-\/-\/-\/-\/-\/-\/-\/-\/-\/-\/-\/-\/-\/-\/-\/-\/-\/-\/-\/-\/-\/-\/-\/-\/-\/-\/-\/-\/-\/-\/-\/-\/-\/-\/-\/-\/-\/-\/-\/-\/-\/-\/-\/-\/-\/-\/-\/-

A leader considers a plan that would seriously harm one individual to
benefit many. The spotlight is Autonomy: meaningful choice and coercion,
plus explicit consent constraints.

Options:

- iphigenia\_coerce Proceed with coerced plan (no valid consent)

- iphigenia\_seek\_consent Seek informed consent with refusal option

- iphigenia\_alternative Choose the alternative that avoids harm

Spotlight coordinates:

iphigenia\_coerce
autonomy\_and\_agency.coercion\_or\_undue\_influence=True;
rights\_and\_duties.has\_valid\_consent=False

iphigenia\_seek\_consent
autonomy\_and\_agency.coercion\_or\_undue\_influence=False;
rights\_and\_duties.has\_valid\_consent=True

iphigenia\_alternative
autonomy\_and\_agency.coercion\_or\_undue\_influence=False;
rights\_and\_duties.has\_valid\_consent=True

-\/-\/- Option: iphigenia\_coerce -\/-\/-

{[}EM=case\_study\_1\_triage {]} verdict=forbid score=0.000

- Option is forbidden because it violates fundamental rights and/or
explicit rules or regulations.

- • violates\_rights = True

- • violates\_explicit\_rule = True

{[}EM=rights\_first\_compliance {]} verdict=forbid score=0.000

- Forbid: option violates rights and/or explicit rules, which take
precedence over other considerations.

- • violates\_rights = True

- • violates\_explicit\_rule = True

{[}EM=tragic\_conflict {]} verdict=avoid score=0.191

- Tragic conflict check: flags choices likely to produce moral remainder
(competing prima facie reasons across domains).

- Tragic conflict index=0.90 (higher means more tension across domains).

- Trigger(s): high\_urgency\_vs\_missing\_procedure,
high\_benefit\_vs\_coercion, rights\_or\_rule\_violation\_present.

- Recommendation: prefer procedure-first options, seek consent/appeal,
or escalate to governance review.

- • tragic\_conflict\_high = True

{[}EM=geneva\_baseline {]} verdict=forbid score=0.000

- Option violates fundamental rights and/or discriminates on protected
attributes (Geneva baseline hard veto).

- • violates\_rights = True

{[}AGG governance{]} verdict=forbid score=0.064

* Marked FORBIDDEN by governance.

* Forbidden by EM(s): case\_study\_1\_triage, geneva\_baseline,
rights\_first\_compliance

* Veto EM(s): geneva\_baseline, rights\_first\_compliance

-\/-\/- Option: iphigenia\_seek\_consent -\/-\/-

{[}EM=case\_study\_1\_triage {]} verdict=neutral score=0.414

- Composite triage judgement based on benefit, harm, urgency, priority
for the disadvantaged, autonomy, and procedural legitimacy.

- Epistemic penalty applied: uncertainty\_level=0.35,
evidence\_quality=medium, novel\_situation\_flag=False. Multiplier=0.82.

{[}EM=rights\_first\_compliance {]} verdict=prefer score=0.800

- Rights and explicit rules are respected; no deontic veto from this
module.

{[}EM=tragic\_conflict {]} verdict=prefer score=0.821

- Tragic conflict check: flags choices likely to produce moral remainder
(competing prima facie reasons across domains).

- Tragic conflict index=0.00 (higher means more tension across domains).

- • tragic\_conflict\_high = False

{[}EM=geneva\_baseline {]} verdict=strongly\_prefer score=0.827

- Epistemic adjustment: multiplier=0.88 (penalty=0.12).

{[}AGG governance{]} verdict=prefer score=0.759

-\/-\/- Option: iphigenia\_alternative -\/-\/-

{[}EM=case\_study\_1\_triage {]} verdict=neutral score=0.519

- Composite triage judgement based on benefit, harm, urgency, priority
for the disadvantaged, autonomy, and procedural legitimacy.

- Epistemic penalty applied: uncertainty\_level=0.20,
evidence\_quality=high, novel\_situation\_flag=False. Multiplier=0.92.

{[}EM=rights\_first\_compliance {]} verdict=prefer score=0.800

- Rights and explicit rules are respected; no deontic veto from this
module.

{[}EM=tragic\_conflict {]} verdict=prefer score=0.885

- Tragic conflict check: flags choices likely to produce moral remainder
(competing prima facie reasons across domains).

- Tragic conflict index=0.00 (higher means more tension across domains).

- • tragic\_conflict\_high = False

{[}EM=geneva\_baseline {]} verdict=strongly\_prefer score=0.902

- Epistemic adjustment: multiplier=0.96 (penalty=0.04).

{[}AGG governance{]} verdict=strongly\_prefer score=0.821

=== Governance Outcome (profile \#1) ===

Selected option:
\textquotesingle iphigenia\_alternative\textquotesingle{}

Ranked options (eligible):
{[}\textquotesingle iphigenia\_alternative\textquotesingle,
\textquotesingle iphigenia\_seek\_consent\textquotesingle{]}

Forbidden options:
{[}\textquotesingle iphigenia\_coerce\textquotesingle{]}

Rationale: Selected option
\textquotesingle iphigenia\_alternative\textquotesingle{} based on
aggregated normative scores and
GovernanceConfig(min\_score\_threshold=0.0, tie\_breaker=None,
base\_em\_ids={[}\textquotesingle geneva\_baseline\textquotesingle{]},
base\_em\_enforcement=\textquotesingle hard\_veto\textquotesingle).
Forbidden options:
{[}\textquotesingle iphigenia\_coerce\textquotesingle{]}.

======================================================================================

=== Pantheon Case 5/8: Hippolytus: The Private Letter ===

Spotlight domain: Privacy \& Data Governance

-\/-\/-\/-\/-\/-\/-\/-\/-\/-\/-\/-\/-\/-\/-\/-\/-\/-\/-\/-\/-\/-\/-\/-\/-\/-\/-\/-\/-\/-\/-\/-\/-\/-\/-\/-\/-\/-\/-\/-\/-\/-\/-\/-\/-\/-\/-\/-\/-\/-\/-\/-\/-\/-\/-\/-\/-\/-\/-\/-\/-\/-\/-\/-\/-\/-\/-\/-\/-\/-\/-\/-\/-\/-\/-\/-\/-\/-\/-\/-\/-\/-\/-\/-\/-\/-

A private accusation letter exists. Options include publishing it
widely, sharing it only with a neutral adjudicator, or destroying it
(protecting privacy but undermining accountability). Spotlight: Privacy
\& Data Governance.

Options:

- hippolytus\_publish Publish the letter publicly (maximum exposure)

- hippolytus\_confidential Share with neutral adjudicator (minimize
disclosure)

- hippolytus\_destroy Destroy the letter (privacy-maximal,
accountability-minimal)

Spotlight coordinates:

hippolytus\_publish privacy\_and\_data.privacy\_invasion\_level=0.95;
privacy\_and\_data.secondary\_use\_without\_consent=True

hippolytus\_confidential
privacy\_and\_data.privacy\_invasion\_level=0.2;
privacy\_and\_data.secondary\_use\_without\_consent=False

hippolytus\_destroy privacy\_and\_data.privacy\_invasion\_level=0.05;
privacy\_and\_data.secondary\_use\_without\_consent=False

-\/-\/- Option: hippolytus\_publish -\/-\/-

{[}EM=case\_study\_1\_triage {]} verdict=forbid score=0.000

- Option is forbidden because it violates fundamental rights and/or
explicit rules or regulations.

- • violates\_rights = True

provenance: rights\_and\_duties.violates\_rights ← hybrid
(rule\_id=RIGHTS-PRIV-010, conf=0.85)

notes: Derived: severe privacy invasion treated as rights violation.

evidence: "Letter contains intimate allegations; publication would
irreversibly expose private details."

- • violates\_explicit\_rule = True

{[}EM=rights\_first\_compliance {]} verdict=forbid score=0.000

- Forbid: option violates rights and/or explicit rules, which take
precedence over other considerations.

- • violates\_rights = True

provenance: rights\_and\_duties.violates\_rights ← hybrid
(rule\_id=RIGHTS-PRIV-010, conf=0.85)

notes: Derived: severe privacy invasion treated as rights violation.

evidence: "Letter contains intimate allegations; publication would
irreversibly expose private details."

- • violates\_explicit\_rule = True

{[}EM=tragic\_conflict {]} verdict=neutral score=0.445

- Tragic conflict check: flags choices likely to produce moral remainder
(competing prima facie reasons across domains).

- Tragic conflict index=0.55 (higher means more tension across domains).

- Trigger(s): severe\_privacy\_vs\_accountability,
rights\_or\_rule\_violation\_present.

- Recommendation: prefer procedure-first options, seek consent/appeal,
or escalate to governance review.

- • tragic\_conflict\_high = True

{[}EM=geneva\_baseline {]} verdict=forbid score=0.000

- Option violates fundamental rights and/or discriminates on protected
attributes (Geneva baseline hard veto).

- • violates\_rights = True

provenance: rights\_and\_duties.violates\_rights ← hybrid
(rule\_id=RIGHTS-PRIV-010, conf=0.85)

notes: Derived: severe privacy invasion treated as rights violation.

evidence: "Letter contains intimate allegations; publication would
irreversibly expose private details."

{[}AGG governance{]} verdict=forbid score=0.148

* Marked FORBIDDEN by governance.

* Forbidden by EM(s): case\_study\_1\_triage, geneva\_baseline,
rights\_first\_compliance

* Veto EM(s): geneva\_baseline, rights\_first\_compliance

-\/-\/- Option: hippolytus\_confidential -\/-\/-

{[}EM=case\_study\_1\_triage {]} verdict=neutral score=0.511

- Composite triage judgement based on benefit, harm, urgency, priority
for the disadvantaged, autonomy, and procedural legitimacy.

- Epistemic penalty applied: uncertainty\_level=0.30,
evidence\_quality=medium, novel\_situation\_flag=False. Multiplier=0.84.

{[}EM=rights\_first\_compliance {]} verdict=prefer score=0.800

- Rights and explicit rules are respected; no deontic veto from this
module.

{[}EM=tragic\_conflict {]} verdict=prefer score=0.875

- Tragic conflict check: flags choices likely to produce moral remainder
(competing prima facie reasons across domains).

- Tragic conflict index=0.00 (higher means more tension across domains).

- • tragic\_conflict\_high = False

{[}EM=geneva\_baseline {]} verdict=strongly\_prefer score=0.810

- Epistemic adjustment: multiplier=0.89 (penalty=0.11).

{[}AGG governance{]} verdict=prefer score=0.785

-\/-\/- Option: hippolytus\_destroy -\/-\/-

{[}EM=case\_study\_1\_triage {]} verdict=avoid score=0.327

- Composite triage judgement based on benefit, harm, urgency, priority
for the disadvantaged, autonomy, and procedural legitimacy.

- Epistemic penalty applied: uncertainty\_level=0.25,
evidence\_quality=medium, novel\_situation\_flag=False. Multiplier=0.85.

{[}EM=rights\_first\_compliance {]} verdict=prefer score=0.800

- Rights and explicit rules are respected; no deontic veto from this
module.

{[}EM=tragic\_conflict {]} verdict=prefer score=0.880

- Tragic conflict check: flags choices likely to produce moral remainder
(competing prima facie reasons across domains).

- Tragic conflict index=0.00 (higher means more tension across domains).

- • tragic\_conflict\_high = False

{[}EM=geneva\_baseline {]} verdict=neutral score=0.553

- Did not follow approved procedure.

- Stakeholders not meaningfully consulted.

- Decision not explainable to the public.

- No meaningful contestation / appeal path.

- Epistemic adjustment: multiplier=0.90 (penalty=0.10).

{[}AGG governance{]} verdict=prefer score=0.674

=== Governance Outcome (profile \#1) ===

Selected option:
\textquotesingle hippolytus\_confidential\textquotesingle{}

Ranked options (eligible):
{[}\textquotesingle hippolytus\_confidential\textquotesingle,
\textquotesingle hippolytus\_destroy\textquotesingle{]}

Forbidden options:
{[}\textquotesingle hippolytus\_publish\textquotesingle{]}

Rationale: Selected option
\textquotesingle hippolytus\_confidential\textquotesingle{} based on
aggregated normative scores and
GovernanceConfig(min\_score\_threshold=0.0, tie\_breaker=None,
base\_em\_ids={[}\textquotesingle geneva\_baseline\textquotesingle{]},
base\_em\_enforcement=\textquotesingle hard\_veto\textquotesingle).
Forbidden options:
{[}\textquotesingle hippolytus\_publish\textquotesingle{]}.

======================================================================================

=== Pantheon Case 6/8: Prometheus: Release of Fire ===

Spotlight domain: Societal \& Environmental

-\/-\/-\/-\/-\/-\/-\/-\/-\/-\/-\/-\/-\/-\/-\/-\/-\/-\/-\/-\/-\/-\/-\/-\/-\/-\/-\/-\/-\/-\/-\/-\/-\/-\/-\/-\/-\/-\/-\/-\/-\/-\/-\/-\/-\/-\/-\/-\/-\/-\/-\/-\/-\/-\/-\/-\/-\/-\/-\/-\/-\/-\/-\/-\/-\/-\/-\/-\/-\/-\/-\/-\/-\/-\/-\/-\/-\/-\/-\/-\/-\/-\/-\/-\/-\/-

A powerful capability (\textquotesingle fire\textquotesingle) can be
released. Options include unrestricted release, release with safety
constraints/training, or withholding. Spotlight: Societal \&
Environmental impacts.

Options:

- prometheus\_unrestricted Unrestricted release (high benefit, high
long-term risk)

- prometheus\_governed Governed release (constraints + oversight)

- prometheus\_withhold Withhold capability (low risk, low benefit)

Spotlight coordinates:

prometheus\_unrestricted
societal\_and\_environmental.long\_term\_societal\_risk=0.85;
societal\_and\_environmental.environmental\_harm=0.55

prometheus\_governed
societal\_and\_environmental.long\_term\_societal\_risk=0.45;
societal\_and\_environmental.environmental\_harm=0.25

prometheus\_withhold
societal\_and\_environmental.long\_term\_societal\_risk=0.15;
societal\_and\_environmental.environmental\_harm=0.05

-\/-\/- Option: prometheus\_unrestricted -\/-\/-

{[}EM=case\_study\_1\_triage {]} verdict=neutral score=0.433

- Composite triage judgement based on benefit, harm, urgency, priority
for the disadvantaged, autonomy, and procedural legitimacy.

- Epistemic penalty applied: uncertainty\_level=0.40,
evidence\_quality=medium, novel\_situation\_flag=False. Multiplier=0.80.

{[}EM=rights\_first\_compliance {]} verdict=prefer score=0.800

- Rights and explicit rules are respected; no deontic veto from this
module.

{[}EM=tragic\_conflict {]} verdict=prefer score=0.845

- Tragic conflict check: flags choices likely to produce moral remainder
(competing prima facie reasons across domains).

- Tragic conflict index=0.00 (higher means more tension across domains).

- • tragic\_conflict\_high = False

{[}EM=geneva\_baseline {]} verdict=avoid score=0.322

- High burden on vulnerable groups.

- Did not follow approved procedure.

- Stakeholders not meaningfully consulted.

- Decision not explainable to the public.

- No meaningful contestation / appeal path.

- Epistemic adjustment: multiplier=0.87 (penalty=0.13).

{[}AGG governance{]} verdict=prefer score=0.601

-\/-\/- Option: prometheus\_governed -\/-\/-

{[}EM=case\_study\_1\_triage {]} verdict=prefer score=0.604

- Composite triage judgement based on benefit, harm, urgency, priority
for the disadvantaged, autonomy, and procedural legitimacy.

- Epistemic penalty applied: uncertainty\_level=0.25,
evidence\_quality=high, novel\_situation\_flag=False. Multiplier=0.90.

{[}EM=rights\_first\_compliance {]} verdict=prefer score=0.800

- Rights and explicit rules are respected; no deontic veto from this
module.

{[}EM=tragic\_conflict {]} verdict=prefer score=0.870

- Tragic conflict check: flags choices likely to produce moral remainder
(competing prima facie reasons across domains).

- Tragic conflict index=0.00 (higher means more tension across domains).

- • tragic\_conflict\_high = False

{[}EM=geneva\_baseline {]} verdict=strongly\_prefer score=0.898

- Epistemic adjustment: multiplier=0.95 (penalty=0.05).

{[}AGG governance{]} verdict=strongly\_prefer score=0.827

-\/-\/- Option: prometheus\_withhold -\/-\/-

{[}EM=case\_study\_1\_triage {]} verdict=neutral score=0.439

- Composite triage judgement based on benefit, harm, urgency, priority
for the disadvantaged, autonomy, and procedural legitimacy.

- Epistemic penalty applied: uncertainty\_level=0.20,
evidence\_quality=high, novel\_situation\_flag=False. Multiplier=0.92.

{[}EM=rights\_first\_compliance {]} verdict=prefer score=0.800

- Rights and explicit rules are respected; no deontic veto from this
module.

{[}EM=tragic\_conflict {]} verdict=prefer score=0.890

- Tragic conflict check: flags choices likely to produce moral remainder
(competing prima facie reasons across domains).

- Tragic conflict index=0.00 (higher means more tension across domains).

- • tragic\_conflict\_high = False

{[}EM=geneva\_baseline {]} verdict=strongly\_prefer score=0.859

- Expected benefit is very low.

- Epistemic adjustment: multiplier=0.96 (penalty=0.04).

{[}AGG governance{]} verdict=prefer score=0.796

=== Governance Outcome (profile \#1) ===

Selected option: \textquotesingle prometheus\_governed\textquotesingle{}

Ranked options (eligible):
{[}\textquotesingle prometheus\_governed\textquotesingle,
\textquotesingle prometheus\_withhold\textquotesingle,
\textquotesingle prometheus\_unrestricted\textquotesingle{]}

Forbidden options: {[}{]}

Rationale: Selected option
\textquotesingle prometheus\_governed\textquotesingle{} based on
aggregated normative scores and
GovernanceConfig(min\_score\_threshold=0.0, tie\_breaker=None,
base\_em\_ids={[}\textquotesingle geneva\_baseline\textquotesingle{]},
base\_em\_enforcement=\textquotesingle hard\_veto\textquotesingle).
Forbidden options: none.

======================================================================================

=== Pantheon Case 7/8: Thebes: Inquiry in Crisis ===

Spotlight domain: Procedural \& Legitimacy

-\/-\/-\/-\/-\/-\/-\/-\/-\/-\/-\/-\/-\/-\/-\/-\/-\/-\/-\/-\/-\/-\/-\/-\/-\/-\/-\/-\/-\/-\/-\/-\/-\/-\/-\/-\/-\/-\/-\/-\/-\/-\/-\/-\/-\/-\/-\/-\/-\/-\/-\/-\/-\/-\/-\/-\/-\/-\/-\/-\/-\/-\/-\/-\/-\/-\/-\/-\/-\/-\/-\/-\/-\/-\/-\/-\/-\/-\/-\/-\/-\/-\/-\/-\/-\/-

A city is in crisis and must investigate a cause. Options include a
transparent public inquiry, a secretive inquiry, or suppressing inquiry
to preserve stability. Spotlight: Procedural \& Legitimacy.

Options:

- thebes\_public\_inquiry Run a transparent public inquiry

- thebes\_secret\_inquiry Run a secret inquiry (privacy-protecting but
opaque)

- thebes\_suppress\_inquiry Suppress inquiry entirely (stability-first)

Spotlight coordinates:

thebes\_public\_inquiry
procedural\_and\_legitimacy.followed\_approved\_procedure=True;
procedural\_and\_legitimacy.contestation\_available=True

thebes\_secret\_inquiry
procedural\_and\_legitimacy.followed\_approved\_procedure=True;
procedural\_and\_legitimacy.contestation\_available=False

thebes\_suppress\_inquiry
procedural\_and\_legitimacy.followed\_approved\_procedure=False;
procedural\_and\_legitimacy.contestation\_available=False

-\/-\/- Option: thebes\_public\_inquiry -\/-\/-

{[}EM=case\_study\_1\_triage {]} verdict=neutral score=0.514

- Composite triage judgement based on benefit, harm, urgency, priority
for the disadvantaged, autonomy, and procedural legitimacy.

- Epistemic penalty applied: uncertainty\_level=0.45,
evidence\_quality=medium, novel\_situation\_flag=False. Multiplier=0.78.

{[}EM=rights\_first\_compliance {]} verdict=prefer score=0.800

- Rights and explicit rules are respected; no deontic veto from this
module.

{[}EM=tragic\_conflict {]} verdict=prefer score=0.870

- Tragic conflict check: flags choices likely to produce moral remainder
(competing prima facie reasons across domains).

- Tragic conflict index=0.00 (higher means more tension across domains).

- • tragic\_conflict\_high = False

{[}EM=geneva\_baseline {]} verdict=prefer score=0.744

- Epistemic adjustment: multiplier=0.86 (penalty=0.14).

{[}AGG governance{]} verdict=prefer score=0.762

-\/-\/- Option: thebes\_secret\_inquiry -\/-\/-

{[}EM=case\_study\_1\_triage {]} verdict=avoid score=0.370

- Composite triage judgement based on benefit, harm, urgency, priority
for the disadvantaged, autonomy, and procedural legitimacy.

- Epistemic penalty applied: uncertainty\_level=0.55,
evidence\_quality=low, novel\_situation\_flag=False. Multiplier=0.66.

{[}EM=rights\_first\_compliance {]} verdict=prefer score=0.800

- Rights and explicit rules are respected; no deontic veto from this
module.

{[}EM=tragic\_conflict {]} verdict=prefer score=0.770

- Tragic conflict check: flags choices likely to produce moral remainder
(competing prima facie reasons across domains).

- Tragic conflict index=0.15 (higher means more tension across domains).

- Trigger(s): high\_epistemic\_uncertainty.

- • tragic\_conflict\_high = False

{[}EM=geneva\_baseline {]} verdict=neutral score=0.536

- Stakeholders not meaningfully consulted.

- Decision not explainable to the public.

- No meaningful contestation / appeal path.

- Low evidence quality.

- Epistemic adjustment: multiplier=0.74 (penalty=0.26).

{[}AGG governance{]} verdict=prefer score=0.638

-\/-\/- Option: thebes\_suppress\_inquiry -\/-\/-

{[}EM=case\_study\_1\_triage {]} verdict=avoid score=0.360

- Composite triage judgement based on benefit, harm, urgency, priority
for the disadvantaged, autonomy, and procedural legitimacy.

- Epistemic penalty applied: uncertainty\_level=0.35,
evidence\_quality=medium, novel\_situation\_flag=False. Multiplier=0.82.

{[}EM=rights\_first\_compliance {]} verdict=prefer score=0.800

- Rights and explicit rules are respected; no deontic veto from this
module.

{[}EM=tragic\_conflict {]} verdict=prefer score=0.685

- Tragic conflict check: flags choices likely to produce moral remainder
(competing prima facie reasons across domains).

- Tragic conflict index=0.25 (higher means more tension across domains).

- Trigger(s): high\_urgency\_vs\_missing\_procedure.

- • tragic\_conflict\_high = False

{[}EM=geneva\_baseline {]} verdict=neutral score=0.519

- Did not follow approved procedure.

- Stakeholders not meaningfully consulted.

- Decision not explainable to the public.

- No meaningful contestation / appeal path.

- Epistemic adjustment: multiplier=0.88 (penalty=0.12).

{[}AGG governance{]} verdict=prefer score=0.602

=== Governance Outcome (profile \#1) ===

Selected option:
\textquotesingle thebes\_public\_inquiry\textquotesingle{}

Ranked options (eligible):
{[}\textquotesingle thebes\_public\_inquiry\textquotesingle,
\textquotesingle thebes\_secret\_inquiry\textquotesingle,
\textquotesingle thebes\_suppress\_inquiry\textquotesingle{]}

Forbidden options: {[}{]}

Rationale: Selected option
\textquotesingle thebes\_public\_inquiry\textquotesingle{} based on
aggregated normative scores and
GovernanceConfig(min\_score\_threshold=0.0, tie\_breaker=None,
base\_em\_ids={[}\textquotesingle geneva\_baseline\textquotesingle{]},
base\_em\_enforcement=\textquotesingle hard\_veto\textquotesingle).
Forbidden options: none.

======================================================================================

=== Pantheon Case 8/8: Oedipus: Acting Under Uncertainty ===

Spotlight domain: Epistemic Status

-\/-\/-\/-\/-\/-\/-\/-\/-\/-\/-\/-\/-\/-\/-\/-\/-\/-\/-\/-\/-\/-\/-\/-\/-\/-\/-\/-\/-\/-\/-\/-\/-\/-\/-\/-\/-\/-\/-\/-\/-\/-\/-\/-\/-\/-\/-\/-\/-\/-\/-\/-\/-\/-\/-\/-\/-\/-\/-\/-\/-\/-\/-\/-\/-\/-\/-\/-\/-\/-\/-\/-\/-\/-\/-\/-\/-\/-\/-\/-\/-\/-\/-\/-\/-\/-

A leader must act under uncertain evidence. Options include immediate
punishment based on weak accusations, pausing to gather better evidence,
or escalating to an oracle-like authority. Spotlight: Epistemic Status.

Options:

- oedipus\_punish\_now Punish immediately (low evidence, high
uncertainty)

- oedipus\_gather\_evidence Pause and gather evidence (reduce
uncertainty)

- oedipus\_oracle\_escalate Escalate to oracle authority (ambiguous
evidence)

Spotlight coordinates:

oedipus\_punish\_now epistemic\_status.uncertainty\_level=0.85;
epistemic\_status.evidence\_quality=low

oedipus\_gather\_evidence epistemic\_status.uncertainty\_level=0.25;
epistemic\_status.evidence\_quality=high

oedipus\_oracle\_escalate epistemic\_status.uncertainty\_level=0.55;
epistemic\_status.evidence\_quality=medium

-\/-\/- Option: oedipus\_punish\_now -\/-\/-

{[}EM=case\_study\_1\_triage {]} verdict=forbid score=0.000

- Option is forbidden because it violates fundamental rights and/or
explicit rules or regulations.

- • violates\_rights = True

- • violates\_explicit\_rule = True

{[}EM=rights\_first\_compliance {]} verdict=forbid score=0.000

- Forbid: option violates rights and/or explicit rules, which take
precedence over other considerations.

- • violates\_rights = True

- • violates\_explicit\_rule = True

{[}EM=tragic\_conflict {]} verdict=neutral score=0.310

- Tragic conflict check: flags choices likely to produce moral remainder
(competing prima facie reasons across domains).

- Tragic conflict index=0.75 (higher means more tension across domains).

- Trigger(s): high\_urgency\_vs\_missing\_procedure,
rights\_or\_rule\_violation\_present, high\_epistemic\_uncertainty.

- Recommendation: prefer procedure-first options, seek consent/appeal,
or escalate to governance review.

- • tragic\_conflict\_high = True

{[}EM=geneva\_baseline {]} verdict=forbid score=0.000

- Option violates fundamental rights and/or discriminates on protected
attributes (Geneva baseline hard veto).

- • violates\_rights = True

{[}AGG governance{]} verdict=forbid score=0.103

* Marked FORBIDDEN by governance.

* Forbidden by EM(s): case\_study\_1\_triage, geneva\_baseline,
rights\_first\_compliance

* Veto EM(s): geneva\_baseline, rights\_first\_compliance

-\/-\/- Option: oedipus\_gather\_evidence -\/-\/-

{[}EM=case\_study\_1\_triage {]} verdict=neutral score=0.576

- Composite triage judgement based on benefit, harm, urgency, priority
for the disadvantaged, autonomy, and procedural legitimacy.

- Epistemic penalty applied: uncertainty\_level=0.25,
evidence\_quality=high, novel\_situation\_flag=False. Multiplier=0.90.

{[}EM=rights\_first\_compliance {]} verdict=prefer score=0.800

- Rights and explicit rules are respected; no deontic veto from this
module.

{[}EM=tragic\_conflict {]} verdict=prefer score=0.875

- Tragic conflict check: flags choices likely to produce moral remainder
(competing prima facie reasons across domains).

- Tragic conflict index=0.00 (higher means more tension across domains).

- • tragic\_conflict\_high = False

{[}EM=geneva\_baseline {]} verdict=strongly\_prefer score=0.893

- Epistemic adjustment: multiplier=0.95 (penalty=0.05).

{[}AGG governance{]} verdict=strongly\_prefer score=0.823

-\/-\/- Option: oedipus\_oracle\_escalate -\/-\/-

{[}EM=case\_study\_1\_triage {]} verdict=avoid score=0.377

- Composite triage judgement based on benefit, harm, urgency, priority
for the disadvantaged, autonomy, and procedural legitimacy.

- Epistemic penalty applied: uncertainty\_level=0.55,
evidence\_quality=medium, novel\_situation\_flag=True. Multiplier=0.67.

{[}EM=rights\_first\_compliance {]} verdict=prefer score=0.800

- Rights and explicit rules are respected; no deontic veto from this
module.

{[}EM=tragic\_conflict {]} verdict=prefer score=0.865

- Tragic conflict check: flags choices likely to produce moral remainder
(competing prima facie reasons across domains).

- Tragic conflict index=0.00 (higher means more tension across domains).

- • tragic\_conflict\_high = False

{[}EM=geneva\_baseline {]} verdict=neutral score=0.511

- Stakeholders not meaningfully consulted.

- Decision not explainable to the public.

- No meaningful contestation / appeal path.

- Novel situation → Geneva baseline is more cautious.

- Epistemic adjustment: multiplier=0.69 (penalty=0.31).

{[}AGG governance{]} verdict=prefer score=0.662

=== Governance Outcome (profile \#1) ===

Selected option:
\textquotesingle oedipus\_gather\_evidence\textquotesingle{}

Ranked options (eligible):
{[}\textquotesingle oedipus\_gather\_evidence\textquotesingle,
\textquotesingle oedipus\_oracle\_escalate\textquotesingle{]}

Forbidden options:
{[}\textquotesingle oedipus\_punish\_now\textquotesingle{]}

Rationale: Selected option
\textquotesingle oedipus\_gather\_evidence\textquotesingle{} based on
aggregated normative scores and
GovernanceConfig(min\_score\_threshold=0.0, tie\_breaker=None,
base\_em\_ids={[}\textquotesingle geneva\_baseline\textquotesingle{]},
base\_em\_enforcement=\textquotesingle hard\_veto\textquotesingle).
Forbidden options:
{[}\textquotesingle oedipus\_punish\_now\textquotesingle{]}.

======================================================================================

Pantheon summary (profile \#1):

- aulis {[}Consequences {]} -\textgreater{} aulis\_delay

- antigone {[}Rights \& Duties {]} -\textgreater{} antigone\_petition

- ajax {[}Justice \& Fairness {]} -\textgreater{} ajax\_merit\_award

- iphigenia {[}Autonomy \& Agency {]} -\textgreater{}
iphigenia\_alternative

- hippolytus {[}Privacy \& Data Governance{]} -\textgreater{}
hippolytus\_confidential

- prometheus {[}Societal \& Environmental{]} -\textgreater{}
prometheus\_governed

- thebes {[}Procedural \& Legitimacy {]} -\textgreater{}
thebes\_public\_inquiry

- oedipus {[}Epistemic Status {]} -\textgreater{}
oedipus\_gather\_evidence
\end{quote}

\end{document}
